% block: nonhermitian-systems-heading page:111 type:section
\subsection{{Non-Hermitian Linear Systems}}\label{sec:nonhermitian-linear-systems}
{Let us recall that the CG method and the CR method (or the MINRES method) have the following two favorable properties:}
\begin{enumerate}
\item{The residual vectors satisfy Ritz-Galerkin approach \eqref{eq:ritz-galerkin-orthogonality} or minimal residual approach \eqref{eq:minimal-residual-problem}.}
\item{The residual vectors are generated by short-term recurrences.}
\end{enumerate}
{As described in Sect.\ \ref{sec:hermitian-linear-systems}, the first property determines the optimality of approximate solutions, and the second property plays an important role in keeping the computational costs and memory requirements constant at each iteration step. Recalling that the Ritz--Galerkin approach \eqref{eq:ritz-galerkin-orthogonality} is for the CG method, and the minimal residual approach \eqref{eq:minimal-residual-problem} is for the CR method and the MINRES method.}
% block: faber-manteuffel-limitation page:112 type:ordinary-paragraph
{Here, a question arises: are there Krylov subspace methods satisfying both the properties for non-Hermitian linear systems? The answer is negative, which was proved in \BookCite{57} and known as the Faber--Manteuffel theorem. More concretely, if matrix $A$ has the form}
\[A=e^{i\phi}(icI+G)\quad(c\ge0,\ 0\le\phi\le2\pi,\ B^{-1}G^{\mathrm H}B=G),\]
{then algorithms satisfying both properties can be constructed. For details, see \BookCite{57} and Greenbaum's book \BookCite{81}, Chap.6. See also \BookCite{121} and the references therein.}
% block: nonhermitian-normal-equations pages:112,113 type:ordinary-paragraph
{For solving non-Hermitian linear systems, one may consider the following transformed linear systems:}
\begin{equation}\label{eq:nonhermitian-normal-equations}A^{\mathrm H}A\bm x=A^{\mathrm H}\bm b.\end{equation}
{If $A$ is nonsingular, then $A^{\mathrm H}A$ is Hermitian positive definite. Thus the CG method can be applied to the transformed linear systems. However, the condition number of $A^{\mathrm H}A$ may be twice as large as the condition number of the original matrix $A$. Hence, from Theorem\ \ref{thm:cg-condition-number-bound}, the speed of convergence may be slow. Furthermore, if the condition number is very large, then it follows from the error analysis \eqref{eq:relative-perturbation-bound}, the accuracy of an obtained approximate solution may be very low. Note that the error analysis \eqref{eq:relative-perturbation-bound} does not depend on numerical algorithms for solving linear systems. This means that the accuracy of an approximate solution computed by any numerical algorithm may become low.}
{On the other hand, if the condition number $A^{\mathrm H}A$ is not so large, there are cases where applying the CG method to \eqref{eq:nonhermitian-normal-equations} is a method of choice. In fact, if $A$ is a unitary matrix, then $A^{\mathrm H}A=I$, which indicates that the CG method gives the solution within only one iteration. Thus, we can guess that this approach is particularly useful for the case where $A$ is close to a unitary matrix. Well-known algorithms in this approach are the CGNE method and the CGNR method, which will be described in Sect.\ \ref{subsec:cg-normal-equations}.}
{In what follows, Krylov subspace methods for non-Hermitian linear systems are derived. We will see that the BiCG method, the BiCR method, and the QMR method can be regarded as extensions of the CG method, the CR method, and the MINRES method respectively, and these are based on the bi-Lanczos process (or a bi-Lanczos-like process), and the GMRES method can be regarded as an extension of the MINRES method, and the GMRES method is based on the Arnoldi process. We will also see how the BiCG method has been improved via product-type Krylov subspace methods and the IDR($s$) method.}
% block: bicg-heading page:113 type:subsection
\subsubsection{{The Bi-Conjugate Gradient (BiCG) Method}}\label{subsec:bicg-method}
{Similar to the derivation process of the CG method in Sect.\ \ref{subsec:conjugate-gradient}, the BiCG method \BookCite{58}, \BookCite{119} can be derived by the bi-Lanczos process in Section\ \ref{subsec:bi-lanczos-process}.}
{In this subsection, we give a concise way of the derivation to see that the BiCG method for non-Hermitian linear systems can be derived from the preconditioned COCG method (Algorithm\ \ref{alg:pcocg}). First, we consider the following $2N\times2N$ complex symmetric linear system:}
\begin{equation}\label{eq:bicg-augmented-system}
\begin{bmatrix}O&A\\A^{\top}&O\end{bmatrix}\begin{bmatrix}\bar{\bm x}^{*}\\\bm x\end{bmatrix}=\begin{bmatrix}\bm b\\\bar{\bm b}^{*}\end{bmatrix},\quad\text{or }\widetilde A\widetilde{\bm x}=\widetilde{\bm b}.
\end{equation}
{We can see from the system \eqref{eq:bicg-augmented-system} that it is mathematically equivalent to}
\[A\bm x=\bm b\quad\text{and}\quad A^{\top}\bar{\bm x}^{*}=\bar{\bm b}^{*}\quad(\Leftrightarrow A^{\mathrm H}\bm x^{*}=\bm b^{*}).\]
{If we apply the algorithm of COCG with the preconditioner}
\begin{equation}\label{eq:bicg-swap-preconditioner}\widetilde M=\begin{bmatrix}O&I\\I&O\end{bmatrix},\quad I:\text{ identity matrix},\end{equation}
{to \eqref{eq:bicg-augmented-system}, then the resulting algorithm at the $n$th iteration step can be written as}
% block: bicg-augmented-cocg-recurrences page:114 type:ordinary-paragraph
\[\begin{aligned}
\widetilde{\bm p}_n^{\mathrm{COCG}}&=\widetilde M^{-1}\widetilde{\bm r}_n^{\mathrm{COCG}}+\beta_{n-1}\widetilde{\bm p}_{n-1}^{\mathrm{COCG}},\\
\alpha_n&=\frac{(\overline{\widetilde M^{-1}\widetilde{\bm r}_n^{\mathrm{COCG}}},\widetilde{\bm r}_n^{\mathrm{COCG}})}{(\overline{\widetilde{\bm p}_n^{\mathrm{COCG}}},\widetilde A\widetilde{\bm p}_n^{\mathrm{COCG}})},\\
\widetilde{\bm x}_{n+1}^{\mathrm{COCG}}&=\widetilde{\bm x}_n^{\mathrm{COCG}}+\alpha_n\widetilde{\bm p}_n^{\mathrm{COCG}},\\
\widetilde{\bm r}_{n+1}^{\mathrm{COCG}}&=\widetilde{\bm r}_n^{\mathrm{COCG}}-\alpha_n\widetilde A\widetilde{\bm p}_n^{\mathrm{COCG}},\\
\beta_n&=\frac{(\overline{\widetilde M^{-1}\widetilde{\bm r}_{n+1}^{\mathrm{COCG}}},\widetilde{\bm r}_{n+1}^{\mathrm{COCG}})}{(\overline{\widetilde M^{-1}\widetilde{\bm r}_n^{\mathrm{COCG}}},\widetilde{\bm r}_n^{\mathrm{COCG}})}.
\end{aligned}\]
{Substituting $\widetilde M^{-1}(=\widetilde M)$ of \eqref{eq:bicg-swap-preconditioner} and the vectors}
\[\widetilde{\bm x}_n^{\mathrm{COCG}}:=\begin{bmatrix}\bm x_n\\\bar{\bm x}_n^{*}\end{bmatrix},\quad\widetilde{\bm r}_n^{\mathrm{COCG}}:=\begin{bmatrix}\bm r_n\\\bar{\bm r}_n^{*}\end{bmatrix},\quad\widetilde{\bm p}_n^{\mathrm{COCG}}:=\begin{bmatrix}\bar{\bm p}_n^{*}\\\bm p_n\end{bmatrix}\]
{into the previous recurrences, and using the following results:}
\[\begin{aligned}
(\overline{\widetilde M^{-1}\widetilde{\bm r}_n^{\mathrm{COCG}}},\widetilde{\bm r}_n^{\mathrm{COCG}})&=\begin{bmatrix}\bar{\bm r}_n^{*\top}&\bm r_n^{\top}\end{bmatrix}\begin{bmatrix}\bm r_n\\\bar{\bm r}_n^{*}\end{bmatrix}\\
&=\bar{\bm r}_n^{*\top}\bm r_n+\bm r_n^{\top}\bar{\bm r}_n^{*}\\
&=\bm r_n^{*\mathrm H}\bm r_n+(\bm r_n^{\top}\bar{\bm r}_n^{*})^{\top}\\
&=(\bm r_n^{*},\bm r_n)+\bm r_n^{*\mathrm H}\bm r_n\\
&=2(\bm r_n^{*},\bm r_n)\\
(\overline{\widetilde{\bm p}_n^{\mathrm{COCG}}},\widetilde A\widetilde{\bm p}_n^{\mathrm{COCG}}),&=\begin{bmatrix}((\bar{\bm p}_n^{*})^{\top}\bm p_n)^{\top}\end{bmatrix}\begin{bmatrix}A\bm p_n\\A^{\top}\bar{\bm p}_n^{*}\end{bmatrix}\\
&=(\bm p_n^{*})^{\mathrm H}A\bm p_n+\bm p_n^{\top}A^{\top}\bar{\bm p}_n^{*},\\
&=(\bm p_n^{*})^{\mathrm H}A\bm p_n+(\bm p_n^{\top}A^{\top}\bar{\bm p}_n^{*})^{\top},\\
&=2(\bm p_n^{*},A\bm p_n),
\end{aligned}\]
{we readily obtain the algorithm of BiCG for solving non-Hermitian linear systems described in Algorithm\ \ref{alg:bicg}. If matrix $A$ is real non-symmetric, the aforementioned derivation process corresponds to the derivation in van der Vorst's book \BookCite{196}, Chap.7, where the BiCG method for real non-symmetric linear systems is derived from the CG method.}
{It can be seen from Algorithm\ \ref{alg:bicg} that the choice $\bm r_0^{*}=\bm r_0$ reduces to the CG method if the coefficient matrix $A$ is Hermitian.}
% block: bicg-biorthogonality-proposition page:114 type:proposition
{Similar to the CG iterates as shown in Proposition\ \ref{prop:cg-orthogonality}, the BiCG iterates have the following relations:}
\par\Needspace{3\baselineskip}
\begin{Proposition}\label{prop:bicg-biorthogonality}
{Let $\bm p_n,\bm r_n$ be the vectors in the BiCG method. Then,}
\begin{enumerate}\item $(\bm r_i^{*},\bm r_j)=0\quad\text{for }i\ne j$,\item $(\bm p_i^{*},A\bm p_j)=0\quad\text{for }i\ne j$.\end{enumerate}
\end{Proposition}
% block: bicg-algorithm page:115 type:ordinary-paragraph
\begin{algorithm}[H]\caption{{The BiCG method}}\label{alg:bicg}\begin{algorithmic}[1]
\Statex\textbf{Input:} $\bm x_0\in\mathbb C^N,\ \beta_{-1}=0,\ \bm p_{-1}=\bm p_{-1}^{*}=\bm0,\ \bm r_0=\bm b-A\bm x_0$
\Statex\textbf{Input:} {Choose $\bm r_0^{*}\in\mathbb C^N$, e.g., $\bm r_0^{*}=\bm r_0$}
\Statex\textbf{Output:} $\bm x_n$
\For{$n=0,1,2,\ldots$, {until convergence}}
\State $\bm p_n=\bm r_n+\beta_{n-1}\bm p_{n-1},\quad\bm p_n^{*}=\bm r_n^{*}+\bar\beta_{n-1}\bm p_{n-1}^{*}$
\State $\alpha_n=\frac{(\bm r_n^{*},\bm r_n)}{(\bm p_n^{*},A\bm p_n)}$
\State $\bm x_{n+1}=\bm x_n+\alpha_n\bm p_n$
\State $\bm r_{n+1}=\bm r_n-\alpha_nA\bm p_n,\quad\bm r_{n+1}^{*}=\bm r_n^{*}-\bar\alpha_nA^{\mathrm H}\bm p_n^{*}$
\State $\beta_n=\frac{(\bm r_{n+1}^{*},\bm r_{n+1})}{(\bm r_n^{*},\bm r_n)}$
\EndFor\end{algorithmic}\end{algorithm}
% block: bicg-krylov-corollary page:115 type:corollary
{From Proposition\ \ref{prop:bicg-biorthogonality}, the following properties hold true:}
\par\Needspace{3\baselineskip}
\begin{Corollary}\label{cor:bicg-krylov-orthogonality}
{Let $\bm p_n,\bm r_n$ be the vectors in the BiCG method. Then,}
\begin{enumerate}\item $\bm r_n\perp\mathcal K_n(A^{\mathrm H},\bm r_0^{*})$,\item $A\bm p_n\perp\mathcal K_n(A^{\mathrm H},\bm r_0^{*})$.\end{enumerate}
\end{Corollary}
% block: bicg-alternative-alpha page:115 type:ordinary-paragraph
{In what follows, we give the other computational formulas for $\alpha_n$ and $\beta_n$ of the BiCG method, which will be used to derive the product-type Krylov subspace methods in Section\ \ref{subsec:product-type-krylov}.}
{Since $\bm r_n^{*}\in\mathcal K_{n+1}(A^{\mathrm H},\bm r_0^{*})$, the vector $\bm r_n^{*}$ can be written as}
\[\begin{aligned}\bm r_n^{*}&=c_n(A^{\mathrm H})^n\bm r_0^{*}+\underbrace{c_{n-1}(A^{\mathrm H})^{n-1}\bm r_0^{*}+\cdots+c_0\bm r_0^{*}}_{\bm z}\\&=c_n(A^{\mathrm H})^n\bm r_0^{*}+\bm z,\quad\bm z\in\mathcal K_n(A^{\mathrm H},\bm r_0^{*}).\end{aligned}\]
{From the above equation and Corollary\ \ref{cor:bicg-krylov-orthogonality}, it follows that}
\[\begin{aligned}(\bm r_n^{*},\bm r_n)&=(c_n(A^{\mathrm H})^n\bm r_0^{*}+\bm z,\bm r_n)=\bar c_n((A^{\mathrm H})^n\bm r_0^{*},\bm r_n)+(\bm z,\bm r_n)\\&=\bar c_n((A^{\mathrm H})^n\bm r_0^{*},\bm r_n)\\
(\bm r_n^{*},A\bm p_n)&=(c_n(A^{\mathrm H})^n\bm r_0^{*}+\bm z,A\bm p_n)=\bar c_n((A^{\mathrm H})^n\bm r_0^{*},A\bm p_n)+(\bm z,A\bm p_n)\\&=\bar c_n((A^{\mathrm H})^n\bm r_0^{*},A\bm p_n).
\end{aligned}\]
{Thus, $\alpha_n$ of the BiCG method (Algorithm\ \ref{alg:bicg}) can be rewritten as}
\begin{equation}\label{eq:bicg-alpha-shadow-power}\alpha_n=\frac{(\bm r_n^{*},\bm r_n)}{(\bm p_n^{*},A\bm p_n)}=\frac{((A^{\mathrm H})^n\bm r_0^{*},\bm r_n)}{((A^{\mathrm H})^n\bm r_0^{*},A\bm p_n)}.\end{equation}
% block: bicg-alternative-beta pages:115,116 type:ordinary-paragraph
{Similarly, we give the other computational formulas for $\beta_n$. From Corollary\ \ref{cor:bicg-krylov-orthogonality}, two vectors $(A^{\mathrm H})^n\bm r_0^{*}\in\mathcal K_{n+1}(A^{\mathrm H},\bm r_0^{*})$ and $A\bm p_{n+1}$ are orthogonal, and thus}
\[\begin{aligned}0&=((A^{\mathrm H})^n\bm r_0^{*},A\bm p_{n+1})\\&=((A^{\mathrm H})^n\bm r_0^{*},A\bm r_{n+1}+\beta_nA\bm p_n)\\&=((A^{\mathrm H})^n\bm r_0^{*},A\bm r_{n+1})+\beta_n((A^{\mathrm H})^n\bm r_0^{*},A\bm p_n)\\&=((A^{\mathrm H})^{n+1}\bm r_0^{*},\bm r_{n+1})+\beta_n((A^{\mathrm H})^n\bm r_0^{*},A\bm p_n),\end{aligned}\]
{which immediately leads to}
\[\beta_n=-\frac{((A^{\mathrm H})^{n+1}\bm r_0^{*},\bm r_{n+1})}{((A^{\mathrm H})^n\bm r_0^{*},A\bm p_n)}.\]
{Using \eqref{eq:bicg-alpha-shadow-power} together with the above result yields}
\begin{equation}\label{eq:bicg-beta-shadow-power}\beta_n=-\alpha_n\frac{((A^{\mathrm H})^{n+1}\bm r_0^{*},\bm r_{n+1})}{((A^{\mathrm H})^n\bm r_0^{*},\bm r_n)}.\end{equation}
% block: composite-step-bicg-heading page:116 type:subsection
\subsubsection{{The Composite Step Bi-Conjugate Gradient (CSBiCG) Method}}\label{subsec:composite-step-bicg}
{From the algorithm of the BiCG method (Algorithm\ \ref{alg:bicg}), there are two kinds of possible (near) breakdown:}
\begin{enumerate}
\item $\rho_n:=(\bm r_n^{*},\bm r_n)=0$ {for $\bm r_n\ne\bm0$ (Lanczos breakdown),}
\item $\sigma_n:=(\bm p_n^{*},A\bm p_n)=0$ ({pivot breakdown}).
\end{enumerate}
{The Lanczos breakdown can be circumvented by the look-ahead Lanczos process in \BookCite{145}, \BookCite{189}. Bank and Chan focused on the pivot breakdown and proposed the Composite Step Bi-Conjugate Gradient (CSBiCG) method to cure the pivot breakdown \BookCite{16}, \BookCite{17} under the assumption that the Lanczos breakdown does not occur. Below is a brief explanation of the algorithm that is based on \BookCite{17}.}
% block: composite-step-auxiliary-vectors pages:116,117 type:ordinary-paragraph
{If $\sigma_n=0$, then $\bm r_{n+1}=\bm r_n-(\rho_n/\sigma_n)A\bm p_n$ and $\bm r_{n+1}^{*}=\bm r_n^{*}-\overline{(\rho_n/\sigma_n)}A^{\mathrm H}\bm p_n^{*}$ cannot be obtained. Consider auxiliary vectors:}
\[\begin{aligned}\bm z_{n+1}&=\sigma_n\bm r_n-\rho_nA\bm p_n,\\\bm z_{n+1}^{*}&=\bar\sigma_n\bm r_n-\bar\rho_nA^{\mathrm H}\bm p_n^{*}.\end{aligned}\]
{Then, $\bm z_{n+1}$ and $\bm z_{n+1}^{*}$ exist even if $\sigma_n=0$, and $\bm z_{n+1}$ and $\bm z_{n+1}^{*}$ belong to $\mathcal K_{n+2}(A,\bm r_0)$ and $\mathcal K_{n+2}(A^{\mathrm H},\bm r_0^{*})$ respectively. The composite step technique then considers $\bm r_{n+2}$ and $\bm r_{n+2}^{*}$ as follows:}
\[\begin{aligned}\bm r_{n+2}&=\bm r_n-c_1A\bm p_n-c_2A\bm z_{n+1},\\\bm r_{n+2}^{*}&=\bm r_n^{*}-c_1^{*}A^{\mathrm H}\bm p_n^{*}-c_2^{*}A^{\mathrm H}\bm z_{n+1}^{*},\end{aligned}\]
{where $c_1,c_2,c_1^{*},c_2^{*}$ are determined so that}
\[\bm r_{n+2}\perp\operatorname{span}\{\bm p_n^{*},\bm z_{n+1}^{*}\},\qquad\bm r_{n+2}^{*}\perp\operatorname{span}\{\bm p_n,\bm z_{n+1}\},\]
{i.e., $(\bm r_{n+2},\bm p_n^{*})=(\bm r_{n+2},\bm z_{n+1}^{*})=(\bm r_{n+2}^{*},\bm p_n)=(\bm r_{n+2}^{*},\bm z_{n+1})=0$. Similarly, let}
\[\begin{aligned}\bm p_{n+2}&=\bm r_{n+2}+d_1\bm p_n+d_2\bm z_{n+1},\\\bm p_{n+2}^{*}&=\bm r_{n+2}^{*}+d_1^{*}\bm p_n^{*}+d_2^{*}\bm z_{n+1}^{*},\end{aligned}\]
{and $d_1,d_2,d_1^{*},d_2^{*}$ are determined so that}
\[A\bm p_{n+2}\perp\operatorname{span}\{\bm p_n^{*},\bm z_{n+1}^{*}\},\qquad A^{\mathrm H}\bm p_{n+2}^{*}\perp\operatorname{span}\{\bm p_n,\bm z_{n+1}\}.\]
{Then, it is shown in \BookCite{17}, Theorem 4.4 that $(\bm r_i^{*},\bm r_j)=(\bm p_i^{*},A\bm p_j)=0$ for $i\ne j$, which means that the composite step residual vector $\bm r_{n+2}$ and $A\bm p_{n+2}$ satisfies}
\[\bm r_{n+2}\perp\mathcal K_{n+2}(A^{\mathrm H},\bm r_0^{*}),\quad A\bm p_{n+2}\perp\mathcal K_{n+2}(A^{\mathrm H},\bm r_0^{*}).\]
{which are the same properties that the BiCG method has, see Corollary\ \ref{cor:bicg-krylov-orthogonality}. The resulting algorithm is described in Algorithm\ \ref{alg:composite-step-bicg}.}
% block: composite-step-choice page:117 type:ordinary-paragraph
{From Algorithm\ \ref{alg:composite-step-bicg}, the number of matrix--vector multiplications is the same as that of the BiCG method. The composite step techniques are used for other Krylov subspace methods \BookCite{35}, \BookCite{113}.}
{At each step in Algorithm\ \ref{alg:composite-step-bicg}, one needs to choose a $1\times1$ step (BiCG step) or a $2\times2$ step (composite step). In \BookCite{17}, the following choice is proposed:}
\begin{algorithmic}[1]
\If{$\|\bm z_{n+1}\|\le\|\bm r_n\||\sigma_n|$}
\State $1\times1$ {step}
\Else
\State $\nu_{n+2}=\|\delta_n\bm r_n-\rho_n^3\zeta_{n+1}\bm q_n-\theta_{n+1}\rho_n^2\bm y_{n+1}\|$
\If{$\nu_{n+2}|\sigma_n|<\|\bm z_{n+1}\||\delta_n|$}
\State $2\times2$ {step}
\Else\State $1\times1$ {step}\EndIf\EndIf
\end{algorithmic}
{This choice leads to not only avoiding (near) pivot breakdowns but also some smoothing of the convergence history in residual norms.}
% block: composite-step-bicg-algorithm page:118 type:ordinary-paragraph
\begin{algorithm}[H]\caption{{The Composite Step BiCG method}}\label{alg:composite-step-bicg}\begin{algorithmic}[1]
\Statex\textbf{Input:} $\bm x_0\in\mathbb C^N,\ \bm r_0=\bm b-A\bm x_0$
\Statex\textbf{Input:} {Choose $\bm r_0^{*}\in\mathbb C^N$, e.g., $\bm r_0^{*}=\bm r_0$}
\Statex\textbf{Output:} $\bm x_n$
\State $\bm p_0=\bm r_0,\quad\bm p_0^{*}=\bm r_0^{*}$
\State $\bm q_0=A\bm p_0,\quad\bm q_0^{*}=A^{\mathrm H}\bm p_0^{*}$
\State $\rho_0=(\bm r_0^{*},\bm r_0)$
\State $n=0$
\While{{until convergence}}
\State $\sigma_n=(\bm p_n^{*},\bm q_n)$
\State $\bm z_{n+1}=\sigma_n\bm r_n-\rho_n\bm q_n,\quad\bm z_{n+1}^{*}=\bar\sigma_n\bm r_n^{*}-\bar\rho_n\bm q_n^{*}$
\State $\bm y_{n+1}=A\bm z_{n+1},\quad\bm y_{n+1}^{*}=A^{\mathrm H}\bm z_{n+1}^{*}$
\State $\theta_{n+1}=(\bm z_{n+1}^{*},\bm z_{n+1})$
\State $\zeta_{n+1}=(\bm z_{n+1}^{*},\bm y_{n+1})$
\If{$1\times1$ {step}}
\State $\alpha=\rho_n/\sigma_n$
\State $\rho_{n+1}=\theta_{n+1}/\sigma_n^2$
\State $\beta_{n+1}=\rho_{n+1}/\rho_n$
\State $\bm x_{n+1}=\bm x_n+\alpha_n\bm p_n$
\State $\bm r_{n+1}=\bm r_n-\alpha_n\bm q_n,\quad\bm r_{n+1}^{*}=\bm r_n^{*}-\bar\alpha_n\bm q_n^{*}$
\State $\bm p_{n+1}=\bm z_{n+1}/\sigma_n+\beta_{n+1}\bm p_n,\quad\bm p_{n+1}^{*}=\bm z_{n+1}^{*}/\bar\sigma_n+\bar\beta_{n+1}\bm p_n^{*}$
\State $\bm q_{n+1}=\bm y_{n+1}/\sigma_n+\beta_{n+1}\bm q_n,\quad\bm q_{n+1}^{*}=\bm y_{n+1}^{*}/\bar\sigma_n+\bar\beta_{n+1}\bm q_n^{*}$
\State $n=n+1$
\EndIf
\If{$2\times2$ {step}}
\State $\delta_n=\sigma_n\zeta_{n+1}\rho_n^2-\theta_{n+1}^2$
\State $\alpha_n=\zeta_{n+1}\rho_n^3/\delta_n,\quad\alpha_{n+1}=\theta_{n+1}\rho_n^2/\delta_n$
\State $\bm x_{n+2}=\bm x_n+\alpha_n\bm p_n+\alpha_{n+1}\bm z_{n+1}$
\State $\bm r_{n+2}=\bm r_n-\alpha_n\bm q_n-\alpha_{n+1}\bm y_{n+1},\quad\bm r_{n+2}^{*}=\bm r_n^{*}-\bar\alpha_n\bm q_n^{*}-\bar\alpha_{n+1}\bm y_{n+1}$
\State $\rho_{n+2}=(\bm r_{n+2}^{*},\bm r_{n+2})$
\State $\beta_{n+1}=\rho_{n+2}/\rho_n,\quad\beta_{n+2}=\rho_{n+2}\sigma_n/\theta_{n+1}$
\State $\bm p_{n+2}=\bm r_{n+2}+\beta_{n+1}\bm p_n+\beta_{n+2}\bm z_{n+1},\quad\bm p_{n+2}^{*}=\bm r_{n+2}^{*}+\bar\beta_{n+1}\bm p_n^{*}+\bar\beta_{n+2}\bm z_{n+1}^{*}$
\State $\bm q_{n+2}=A\bm p_{n+2},\quad\bm q_{n+2}^{*}=A^{\mathrm H}\bm p_{n+2}^{*}$
\State $n=n+2$
\EndIf\EndWhile\end{algorithmic}\end{algorithm}
% block: composite-bicg-best-approximation pages:118,119 type:theorem
{The following theorem shows the best approximation result of the (composite step) BiCG method:}
\par\Needspace{14\baselineskip}
\begin{Theorem}[\BookCite{16}]\label{thm:composite-bicg-best-approximation}
{Consider applying the (composite step) BiCG method to real nonsymmetric linear systems with initial guess $\bm x_0=\bm0$ and $\bm r_0^{*}=\bm r_0$. Suppose that for all $\bm v,\bm w\in\mathbb R^N$, we have}
\[|\bm w^{\top}A\bm v|\le\Gamma\|\bm v\|\|\bm w\|,\]
{where $\Gamma$ is a constant independent of $\bm v$ and $\bm w$. Further, suppose that for those steps in the (composite step) BiCG method in which we compute an approximation $\bm x_n$, we have}
\[\inf_{\substack{\bm v\in\mathcal K_n(A,\bm r_0)\\\|\bm v\|=1}}\ \sup_{\substack{\bm w\in\mathcal K_n(A^{\top},\bm r_0^{*})\\\|\bm w\|\le1}}\bm w^{\top}A\bm v\ge\delta_n\ge\delta>0.\]
{Then}
\[\|\bm x-\bm x_n\|\le(1+\Gamma/\delta)\inf_{\bm v\in\mathcal K_n(A,\bm r_0)}\|\bm x-\bm v\|.\]
\end{Theorem}
{Theorem\ \ref{thm:composite-bicg-best-approximation} is a simplified result of Theorem 4.1 in \BookCite{16}. From Theorem\ \ref{thm:composite-bicg-best-approximation}, we see that the (composite step) BiCG method produces an approximate solution close to the best approximate solution in error norm when $\Gamma/\delta$ is close to 0.}
% block: bicr-heading page:119 type:subsection
\subsubsection{{The Bi-Conjugate Residual (BiCR) Method}}\label{subsec:bicr-method}
{The BiCR method \BookCite{169} is an extension of the CR method to non-Hermitian linear systems.}\footnote{{The method was mentioned in the unpublished paper \BookCite{38} by Chronopulous and Ma in 1989, and Dr. Chronopulous emailed the author about this fact on Apr. 6th, 2020.}}
{In this subsection, we show that the BiCR method can be derived from the preconditioned COCR method in Algorithm\ \ref{alg:pcocr}. When we apply the algorithm of the COCR method with the preconditioner \eqref{eq:bicg-swap-preconditioner} to \eqref{eq:bicg-augmented-system}, the resulting algorithm at the $n$th iteration step can be written as}
\[\begin{aligned}
\widetilde{\bm p}_n^{\mathrm{COCR}}&=\widetilde M^{-1}\widetilde{\bm r}_n^{\mathrm{COCR}}+\beta_{n-1}\widetilde{\bm p}_{n-1}^{\mathrm{COCR}},\\
\alpha_n&=\frac{(\overline{\widetilde M^{-1}\widetilde{\bm r}_n^{\mathrm{COCR}}},\widetilde A\widetilde M^{-1}\widetilde{\bm r}_n^{\mathrm{COCR}})}{(\overline{\widetilde A\widetilde{\bm p}_n^{\mathrm{COCR}}},\widetilde M^{-1}\widetilde A\widetilde{\bm p}_n^{\mathrm{COCR}})},\\
\widetilde{\bm x}_{n+1}^{\mathrm{COCR}}&=\widetilde{\bm x}_n^{\mathrm{COCR}}+\alpha_n\widetilde{\bm p}_n^{\mathrm{COCR}},\\
\widetilde{\bm r}_{n+1}^{\mathrm{COCR}}&=\widetilde{\bm r}_n^{\mathrm{COCR}}-\alpha_n\widetilde A\widetilde{\bm p}_n^{\mathrm{COCR}},\\
\beta_n&=\frac{(\overline{\widetilde M^{-1}\widetilde{\bm r}_{n+1}^{\mathrm{COCR}}},\widetilde A\widetilde M^{-1}\widetilde{\bm r}_{n+1}^{\mathrm{COCR}})}{(\overline{\widetilde M^{-1}\widetilde{\bm r}_n^{\mathrm{COCR}}},\widetilde A\widetilde M^{-1}\widetilde{\bm r}_n^{\mathrm{COCR}})}.
\end{aligned}\]
{Substituting $\widetilde M^{-1}(=\widetilde M)$ of \eqref{eq:bicg-swap-preconditioner} and the vectors}
\[\widetilde{\bm x}_n^{\mathrm{COCR}}:=\begin{bmatrix}\bm x_n\\\bar{\bm x}_n^{*}\end{bmatrix},\quad\widetilde{\bm r}_n^{\mathrm{COCR}}:=\begin{bmatrix}\bm r_n\\\bar{\bm r}_n^{*}\end{bmatrix},\quad\widetilde{\bm p}_n^{\mathrm{COCR}}:=\begin{bmatrix}\bar{\bm p}_n^{*}\\\bm p_n\end{bmatrix}\]
{into the previous recurrences, and using the following results:}
% block: bicr-bilinear-reductions page:120 type:ordinary-paragraph
\[\begin{aligned}
(\overline{\widetilde M^{-1}\widetilde{\bm r}_n^{\mathrm{COCR}}},\widetilde A\widetilde M^{-1}\widetilde{\bm r}_n^{\mathrm{COCR}})&=\begin{bmatrix}\bar{\bm r}_n^{*\top}&\bm r_n^{\top}\end{bmatrix}\begin{bmatrix}A\bm r_n\\A^{\top}\bar{\bm r}_n^{*}\end{bmatrix}\\
&=\bar{\bm r}_n^{*\top}A\bm r_n+\bm r_n^{\top}A^{\top}\bar{\bm r}_n^{*}\\
&=\bm r_n^{*\mathrm H}A\bm r_n+(\bm r_n^{\top}A^{\top}\bar{\bm r}_n^{*})^{\top}\\
&=(\bm r_n^{*},A\bm r_n)+\bm r_n^{*\mathrm H}A\bm r_n\\
&=2(\bm r_n^{*},A\bm r_n)
\end{aligned}\]
{and}
\[\begin{aligned}
(\overline{\widetilde A\widetilde{\bm p}_n^{\mathrm{COCR}}},\widetilde M^{-1}\widetilde A\widetilde{\bm p}_n^{\mathrm{COCR}})&=\begin{bmatrix}(A\bm p_n)^{\top}&(A^{\top}\bar{\bm p}_n^{*})^{\top}\end{bmatrix}\begin{bmatrix}A^{\top}\bar{\bm p}_n^{*}\\A\bm p_n\end{bmatrix}\\
&=(A\bm p_n)^{\top}A^{\top}\bar{\bm p}_n^{*}+(A^{\top}\bar{\bm p}_n^{*})^{\top}A\bm p_n\\
&=(\bm p_n^{\top}A^{\top}A^{\top}\bar{\bm p}_n^{*})^{\top}+(A^{\mathrm H}\bm p_n^{*})^{\mathrm H}A\bm p_n\\
&=(A^{\mathrm H}\bm p_n^{*})^{\mathrm H}A\bm p_n+(A^{\mathrm H}\bm p_n^{*},A\bm p_n)\\
&=2(A^{\mathrm H}\bm p_n^{*},A\bm p_n),
\end{aligned}\]
{we readily obtain the algorithm of the BiCR method for solving non-Hermitian linear systems.}
% block: bicr-algorithm page:120 type:ordinary-paragraph
\begin{algorithm}[H]\caption{{The BiCR method}}\label{alg:bicr}\begin{algorithmic}[1]
\Statex\textbf{Input:} $\bm x_0\in\mathbb C^N,\ \beta_{-1}=0,\ \bm p_{-1}=\bm p_{-1}^{*}=\bm0,\ \bm r_0=\bm b-A\bm x_0$
\Statex\textbf{Input:} {Choose $\bm r_0^{*}\in\mathbb C^N$, e.g., $\bm r_0^{*}=\bm r_0$}
\Statex\textbf{Output:} $\bm x_n$
\For{$n=0,1,2,\ldots$, {until convergence}}
\State $\bm p_n=\bm r_n+\beta_{n-1}\bm p_{n-1},\quad\bm p_n^{*}=\bm r_n^{*}+\bar\beta_{n-1}\bm p_{n-1}^{*}$
\State $A\bm p_n=A\bm r_n+\beta_{n-1}A\bm p_{n-1}$
\State $\alpha_n=\frac{(\bm r_n^{*},A\bm r_n)}{(A^{\mathrm H}\bm p_n^{*},A\bm p_n)}$
\State $\bm x_{n+1}=\bm x_n+\alpha_n\bm p_n$
\State $\bm r_{n+1}=\bm r_n-\alpha_nA\bm p_n,\quad\bm r_{n+1}^{*}=\bm r_n^{*}-\bar\alpha_nA^{\mathrm H}\bm p_n^{*}$
\State $\beta_n=\frac{(\bm r_{n+1}^{*},A\bm r_{n+1})}{(\bm r_n^{*},A\bm r_n)}$
\EndFor\end{algorithmic}\end{algorithm}
% block: bicr-biorthogonality-proposition page:120 type:proposition
{Similar to Proposition\ \ref{prop:bicg-biorthogonality}, BiCR iterates hold the following properties:}
\par\Needspace{3\baselineskip}
\begin{Proposition}\label{prop:bicr-biorthogonality}
{Let $\bm p_n,\bm r_n$ be the vectors in the BiCR method. Then,}
\begin{enumerate}\item $(\bm r_i^{*},A\bm r_j)=0\quad\text{for }i\ne j$,\item $(A^{\mathrm H}\bm p_i^{*},A\bm p_j)=0\quad\text{for }i\ne j$.\end{enumerate}
\end{Proposition}
{If the coefficient matrix $A$ is Hermitian or real symmetric, the choice $\bm r_0^{*}=\bm r_0$ reduces to the CR method that generates optimal approximate solutions in terms of residual 2-norms. Thus, if $A$ is close to Hermitian or real symmetric, it is expected that the BiCR method has smooth convergence behavior and generates near minimal residual solutions.}
% block: bicr-convergence-comparison page:121 type:ordinary-paragraph
{To see this, we consider nonsymmetric linear systems \eqref{eq:three-dimensional-difference-system} with parameters $N=20$, $a_1=a_2=a_3=1$, two different numbers $b_1=b_2=b_3=0.1$ (near symmetric matrix) and $b_1=b_2=b_3=10$ (far symmetric matrix), and $c=700$. The horizontal axis is the number of matrix--vector multiplications. The stopping criterion used for the BiCG method and the BiCR method is $\|\bm r_n\|/\|\bm b\|\le10^{-10}$.}
{The results for $b_1=b_2=b_3=0.1$ (near symmetric matrix) are shown in Fig.\ \ref{fig:bicg-bicr-near-symmetric}. We see from Fig.\ \ref{fig:bicg-bicr-near-symmetric} that the BiCR method shows much smoother convergence behavior than the BiCG method.}
{The results for $b_1=b_2=b_3=0.1$ (far symmetric matrix) are shown in Fig.\ \ref{fig:bicg-bicr-far-symmetric}. From Fig.\ \ref{fig:bicg-bicr-far-symmetric}, the BiCG method and the BiCR method have similar convergence behavior, and the BiCR method does not show smooth convergence behavior.}
% block: bicr-near-symmetric-figure page:121 type:figure
\begin{figure}[H]\centering\includegraphics[width=.87\linewidth]{figures/bicg-bicr-near-symmetric.png}
\caption{{Convergence histories of BiCG and BiCR for $b_1=b_2=b_3=0.1$}}\label{fig:bicg-bicr-near-symmetric}\end{figure}
% block: bicr-far-symmetric-figure page:122 type:figure
\begin{figure}[H]\centering\includegraphics[width=.87\linewidth]{figures/bicg-bicr-far-symmetric.png}
\caption{{Convergence histories of BiCG and BiCR for $b_1=b_2=b_3=10$}}\label{fig:bicg-bicr-far-symmetric}\end{figure}
% block: qmr-heading pages:121,122 type:subsection
\subsubsection{{The Quasi-Minimal Residual (QMR) Method}}\label{subsec:qmr-method}
{In this subsection, we describe the idea of the Quasi-Minimal Residual (QMR) method \BookCite{66}. The derivation process of the QMR method is almost the same as that of the QMR\_SYM method in Sect.\ \ref{subsec:qmr-sym}. The main difference is the use not of the complex symmetric Lanczos process but the bi-Lanczos process. Let $V_n$ be the bi-orthogonal basis of $\mathcal K_n(A,\bm r_0)$ via the bi-Lanczos process given in Algorithm\ \ref{alg:bi-lanczos}. Then, the QMR method finds an approximate solution $\bm x_n$ such that $\bm x_n$ lies in the affine space $\bm x_0+\mathcal K_n(A,\bm r_0)$, i.e.,}
\[\bm x_n=\bm x_0+V_n\bm y_n,\quad\bm y_n\in\mathbb C^n.\]
{The corresponding residual vector is given by}
\[\bm r_n=\bm r_0-AV_n\bm y_n.\]
{From the matrix form of the bi-Lanczos process \eqref{eq:bilanczos-matrix-relation}, it follows that}
\[\bm r_n=\bm r_0-V_{n+1}T_{n+1,n}\bm y_n=V_{n+1}(\beta\bm e_1-T_{n+1,n}\bm y_n),\]
{where $\beta:=\|\bm r_0\|$. The above derivation process is very similar to that of the MINRES method; however, in this case, if we choose $\bm y_n$ such that the norm of the residual is minimized, then it may lead to a large amount of computational costs and memory requirement as described in Sect.\ \ref{subsec:qmr-sym}. Similar to the QMR\_SYM method, the QMR method chooses $\bm y_n$ such that}
\begin{equation}\label{eq:qmr-projected-minimization}\bm y_n:=\operatorname*{arg\,min}_{\bm y\in\mathbb C^n}\|\beta\bm e_1-T_{n+1,n}\bm y\|.\end{equation}
{The minimization problem can be solved efficiently by Givens rotations. Following the solution of \eqref{eq:minres-projected-least-squares} and Sect.\ \ref{subsec:qmr-sym}, the QMR method is obtained and the algorithm is described in Algorithm\ \ref{alg:qmr}.}
% block: qmr-breakdown pages:122,123 type:ordinary-paragraph
{Finally, one should note that Algorithm\ \ref{alg:qmr}, as well as the BiCG method and the BiCR method, has a possibility of breakdown (zero division). In finite precision arithmetic, such breakdown is very rare; however, near breakdown may occur, and this causes numerical instability. Hence, to avoid the (near) breakdown problem, the QMR method in \BookCite{66} uses the look-ahead Lanczos process, which was proposed by Taylor \BookCite{189} and Parlett et al.\ \BookCite{145}.}
% block: qmr-algorithm page:123 type:ordinary-paragraph
\begin{algorithm}[H]\caption{{The QMR method}}\label{alg:qmr}\begin{algorithmic}[1]
\Statex\textbf{Input:} $\bm x_0\in\mathbb C^N,\ \beta_0=0,\ \gamma_0=0,\ \bm v_0=\bm0,\ \bm w_0=\bm0,\ \bm r_0=\bm b-A\bm x_0$
\Statex\textbf{Output:} $\bm x_n$
\State $\bm g=(\|\bm r_0\|,0,\ldots,0)^{\top},\quad\bm v_1=\bm r_0/\|\bm r_0\|$
\For{$n=1,2,\ldots$}
\State ({Bi-Lanczos process})
\State $\alpha_n=(\bm w_n,A\bm v_n),$
\State $\widetilde{\bm v}_{n+1}=A\bm v_n-\alpha_n\bm v_n-\beta_{n-1}\bm v_{n-1}$
\State $\widetilde{\bm w}_{n+1}=A^{\mathrm H}\bm w_n-\bar\alpha_n\bm w_n-\gamma_{n-1}\bm w_{n-1}$
\State $\gamma_n=\|\widetilde{\bm v}_{n+1}\|_2$
\State $\bm v_{n+1}=\widetilde{\bm v}_{n+1}/\gamma_n$
\State $\beta_n=(\widetilde{\bm w}_{n+1},\bm v_{n+1})$
\State $\bm w_{n+1}=\widetilde{\bm w}_{n+1}/\bar\beta_n$
\State (Givens {rotations})
\State \textbf{for} $i=\max\{1,n-2\},\ldots,n-1$ \textbf{do}
\State $\begin{bmatrix}t_{i,n}\\t_{i+1,n}\end{bmatrix}=\begin{bmatrix}c_i&s_i\\-\bar s_i&c_i\end{bmatrix}\begin{bmatrix}t_{i,n}\\t_{i+1,n}\end{bmatrix}$
\State \textbf{end for}
\State $c_n=\frac{|t_{n,n}|}{\sqrt{|t_{n,n}|^2+|t_{n+1,n}|^2}}$
\State $\bar s_n=\frac{t_{n+1,n}}{t_{n,n}}c_n$
\State $t_{n,n}=c_nt_{n,n}+s_nt_{n+1,n}$
\State $t_{n+1,n}=0$
\State $\begin{bmatrix}g_n\\g_{n+1}\end{bmatrix}=\begin{bmatrix}c_n&s_n\\-\bar s_n&c_n\end{bmatrix}\begin{bmatrix}g_n\\0\end{bmatrix}$
\State ({Update} $\bm x_n$)
\State $\bm p_n=(\bm v_n-t_{n-2,n}\bm p_{n-2},$
\State $\qquad-t_{n-1,n}\bm p_{n-1})/t_{n,n}$
\State $\bm x_n=\bm x_{n-1}+g_n\bm p_n$
\State ({Check convergence})
\State $A\bm p_n=(A\bm v_n-t_{n-2,n}A\bm p_{n-2}$
\State $\qquad-t_{n-1,n}A\bm p_{n-1})/t_{n,n}$
\State $\bm r_n=\bm r_{n-1}-g_nA\bm p_n$
\State {if $\|\bm r_n\|/\|\bm b\|\le\epsilon$, then stop}
\EndFor\end{algorithmic}\end{algorithm}
% block: gmres-heading pages:123,124 type:subsection
\subsubsection{{The Generalized Minimal Residual (GMRES) Method}}\label{subsec:gmres-method}
{The derivation process of the GMRES method \BookCite{152} is closely related to the MINRES method in Sect.\ \ref{subsec:minres}. The GMRES method generates $\bm x_n$ that minimizes $\|\bm b-A\bm x_n\|$ over the affine space $\bm x_0+\mathcal K_n(A,\bm r_0)$. This can efficiently be achieved by using the Arnoldi process in Section\ \ref{subsec:arnoldi-process}. Now, we give the derivation process of the GMRES method.}
{Let $V_n$ be an $N$-by-$n$ matrix whose columns are the orthonormal basis vectors of $\mathcal K_n(A,\bm r_0)$ by the Arnoldi process. Then, since $\bm x_n$ lies in the affine space $\bm x_0+\mathcal K_n(A,\bm r_0)$, we have}
\[\bm x_n=\bm x_0+V_n\bm y_n,\quad\bm y_n\in\mathbb C^n.\]
{The corresponding residual vector is written as}
\[\bm r_n=\bm r_0-AV_n\bm y_n.\]
{From the matrix representation of the Arnoldi process in \eqref{eq:arnoldi-matrix-relation} it follows that}
\begin{equation}\label{eq:gmres-projected-residual}\bm r_n=V_{n+1}(\beta\bm e_1-H_{n+1,n}\bm y_n),\end{equation}
{where $\beta:=\|\bm r_0\|$. Since $V_{n+1}^{\mathrm H}V_{n+1}=I_{n+1}$, i.e., the identity matrix, the 2-norm of the residual vector is given by}
\[\|\bm r_n\|=\|V_{n+1}(\beta\bm e_1-H_{n+1,n}\bm y_n)\|=\|\beta\bm e_1-H_{n+1,n}\bm y_n\|.\]
{Hence, the 2-norm of the residual vector can be minimized by choosing $\bm y_n$ as follows:}
\begin{equation}\label{eq:gmres-projected-minimization}\bm y_n:=\operatorname*{arg\,min}_{\bm y\in\mathbb C^n}\|\beta\bm e_1-H_{n+1,n}\bm y\|.\end{equation}
{The vector $\bm y_n$ can be obtained by using Givens rotations. Following the solution of \eqref{eq:minres-projected-least-squares}, the GMRES method is obtained. Now, we describe the algorithm of GMRES in Algorithm\ \ref{alg:gmres}, which is based on the Arnoldi process of modified Gram--Schmidt type (Algorithm\ \ref{alg:arnoldi-modified}).}
% block: gmres-algorithm page:124 type:ordinary-paragraph
\begin{algorithm}[H]\caption{{The GMRES method of modified Gram-Schmidt type}}\label{alg:gmres}\begin{algorithmic}[1]
\Statex\textbf{Input:} $\bm x_0\in\mathbb C^N,\ \bm r_0=\bm b-A\bm x_0$
\Statex\textbf{Output:} $\bm x_n$
\State $\bm g=(\|\bm r_0\|,0,\ldots,0)^{\top},\quad\bm v_1=\bm r_0/\|\bm r_0\|$
\For{$n=1,2,\ldots$}
\State (Arnoldi {process})
\State $\bm t=A\bm v_n$
\For{$i=1,2,\ldots,n$}
\State $h_{i,n}=(\bm v_i,\bm t)$
\State $\bm t=\bm t-h_{i,n}\bm v_i$
\EndFor
\State $h_{n+1,n}=\|\bm t\|$
\State $\bm v_{n+1}=\bm t/h_{n+1,n}$
\State (Givens {rotations})
\For{$i=1,2,\ldots,n-1$}
\State $\begin{bmatrix}h_{i,n}\\h_{i+1,n}\end{bmatrix}=\begin{bmatrix}c_i&s_i\\-\bar s_i&c_i\end{bmatrix}\begin{bmatrix}h_{i,n}\\h_{i+1,n}\end{bmatrix}$
\EndFor
\State $c_n=\frac{|h_{n,n}|}{\sqrt{|h_{n,n}|^2+|h_{n+1,n}|^2}}$
\State $\bar s_n=\frac{h_{n+1,n}}{h_{n,n}}c_n$
\State $h_{n,n}=c_nh_{n,n}+s_nh_{n+1,n}$
\State $h_{n+1,n}=0$
\State $\begin{bmatrix}g_n\\g_{n+1}\end{bmatrix}=\begin{bmatrix}c_n&s_n\\-\bar s_n&c_n\end{bmatrix}\begin{bmatrix}g_n\\0\end{bmatrix}$
\State ({Check convergence})
\If{$|g_{n+1}|/\|\bm b\|\le\epsilon$}
\State $\bm x_n=\bm x_0+V_nH_n^{-1}\bm g$
\EndIf\EndFor\end{algorithmic}\end{algorithm}
% block: gmres-back-substitution page:124 type:ordinary-paragraph
{Note that computing $\bm c:=H_n^{-1}\bm g$ in line 22 of Algorithm\ \ref{alg:gmres} corresponds to solving linear systems of the form $H_n\bm c=\bm g$. After Givens rotations, $H_n$ becomes not a Hessenberg but an upper triangular matrix. Thus the linear systems can easily be solved by the back substitution as described in \eqref{eq:lu-triangular-solves}.}
{In what follows, the convergence analysis of the GMRES method is described.}
% block: gmres-convergence-proposition pages:124,125 type:proposition
\par\Needspace{3\baselineskip}
\begin{Proposition}\label{prop:gmres-diagonalizable-bound}
{Let $A$ be an $N$-by-$N$ diagonalizable matrix, i.e., $A=X\Lambda X^{-1}$ with $\Lambda=\operatorname{diag}\{\lambda_1,\ldots,\lambda_N\}$, and let $\bm r_n$ be the $m$-step GMRES residual vector. Then}
\[\|\bm r_n\|\le\kappa(X)\epsilon_n\|\bm r_0\|,\]
{where $\kappa(X)=\|X\|\|X^{-1}\|$ and}
\[\epsilon^{(n)}=\min_{p\in\mathbb P_n,p(0)=1}\max_{i=1,\ldots,N}|p(\lambda_i)|.\]
{Here, $\mathbb P_m$ denotes the set of all polynomials of degree at most $m$.}
\end{Proposition}
% block: gmres-convergence-proof page:125 type:proof
\begin{Proof}
{The GMRES method finds an approximate solution $\bm x_n$ such that}
\[\min_{\bm x_n\in\bm x_0+\mathcal K_n(A,\bm r_0)}\|\bm r_n\|.\]
{Since $\bm r_n\in\mathcal K_{n+1}(A,\bm r_0)$, the GMRES residual vector can be written in the matrix polynomial form $\bm r_n=(I+c_1A+c_2A^2+\cdots+c_nA^n)\bm r_0$, so that $c_i$'s are determined by minimizing $\|\bm r_n\|$, or equivalently $\bm r_n=p_n^{(\mathrm{opt})}(A)\bm r_0$ with $p_n^{(\mathrm{opt})}(0)=1$, and $p_n^{(\mathrm{opt})}$ is determined by minimizing $\|\bm r_n\|$, i.e.,}
\[\min_{p\in\mathbb P_n,p(0)=1}\|p_n(A)\bm r_0\|.\]
{Therefore, for any $p\in\mathbb P_n$ and $p(0)=1$, the residual 2-norm $\|\bm r_n\|$ can be bounded by}
\[\begin{aligned}\|\bm r_n\|&=\|p_n^{(\mathrm{opt})}(A)\bm r_0\|\\&\le\|p_n(A)\bm r_0\|\\&=\|Xp_n(\Lambda)X^{-1}\bm r_0\|\\&\le\|X\|\|X^{-1}\|\|p_n(\Lambda)\|\|\bm r_0\|\\&=\kappa(X)\max_{i=1,\ldots,N}|p_n(\lambda_i)|\|\bm r_0\|\end{aligned}\]
{Since the above inequality holds for any $p\in\mathbb P_n$ and $p(0)=1$, we have}
\[\begin{aligned}\|\bm r_n\|&\le\kappa(X)\left(\min_{p\in\mathbb P_n,p(0)=1}\max_{i=1,\ldots,N}|p(\lambda_i)|\right)\|\bm r_0\|\\&=\kappa(X)\epsilon_n\|\bm r_0\|,\end{aligned}\]
{which concludes the proof.}
\par\nopagebreak[4]
\end{Proof}
% block: gmres-eigenvalue-discussion page:125 type:ordinary-paragraph
{From Proposition\ \ref{prop:gmres-diagonalizable-bound}, if the condition number of $X$ is small, the value of $\epsilon_n$ depends highly on the speed of convergence. Further, if the condition number of $X$ is small and eigenvalues are well clustered, the GMRES method shows good convergence behavior. If matrix $A$ is symmetric, then the condition number of $X$ is 1. Therefore the convergence of the GMRES method (or the MINRES method) depends only on the distribution of eigenvalues. By using an upper bound of $\epsilon_n$, the convergence behavior can be estimated without using all the eigenvalues. For the details, see an excellent book by Saad \BookCite{151}, pp.206--207.}
% block: restarted-gmres-motivation page:126 type:ordinary-paragraph
{Nice features of the GMRES method are that the GMRES method never suffers from breakdown, and the residual 2-norm decreases monotonically and has the optimality regarding the residual 2-norm, unlike Krylov subspace methods based on the BiCG method and the BiCR method. On the other hand, a drawback of the GMRES method is that the computational costs and memory requirement of the GMRES method increase linearly with the number of iterations. For example, if the number of iterations is 10000, then an $N\times10000$ matrix needs to be stored and the size of the least-squares problem in \eqref{eq:gmres-projected-minimization} is $10001\times10000$. To overcome the drawback, a restarted version of the GMRES method is proposed, which is described in Algorithm\ \ref{alg:gmres-restarted}.}
{The idea of the restarted GMRES method denoted by GMRES($m$) is that we run the GMRES method until the prescribed maximum iteration number $m$ is reached, and then we restart the GMRES method using the initial guess as the approximate solution produced by the previous GMRES method at the $m$th iteration step.}
% block: restarted-gmres-algorithm page:126 type:ordinary-paragraph
\begin{algorithm}[H]\caption{{The GMRES($m$) method}}\label{alg:gmres-restarted}\begin{algorithmic}[1]
\Statex\textbf{Input:} {Initial guess $\bm x_0\in\mathbb C^N$ and restart number $m$}
\Statex\textbf{Output:} $\bm x_m$
\State{Run $m$ iterations of the GMRES method (Algorithm\ \ref{alg:gmres}) with $\bm x_0$ and produce $\bm x_m$.}
\While{{convergence}}
\State{Set $\bm x_0=\bm x_m$.}
\State{Run $m$ iterations of the GMRES method (Algorithm\ \ref{alg:gmres}) with $\bm x_0$ and produce $\bm x_m$.}
\EndWhile\end{algorithmic}\end{algorithm}
% block: gmres-restart-frequency page:126 type:ordinary-paragraph
{Consider the restart frequency $m$ of the GMRES($m$) method as fixed. This means that we apply the $m$-step GMRES method to the linear systems repeatedly until convergence. On the other hand, an unfixed restart frequency means that we run the GMRES($m_1$) and then run the GMRES($m_2$) method, the GMRES($m_3$) method, and so on, until convergence. If $m_1=m_2=\cdots=m$ then this corresponds to the fixed restart frequency.}
{Efficient ways for determining $m_k$ for $k=1,2,\ldots$ are studied in \BookCite{117} and by some authors, e.g., \BookCite{14}, \BookCite{137}, \BookCite{176}, \BookCite{211}.}
% block: lgmres-augmentation pages:126,127 type:ordinary-paragraph
{Another important development of the GMRES method is augmentation. There are two approaches to the augmentation: one is to use previous approximate solutions, e.g., the LGMRES method \BookCite{15}; another is to use approximate eigenvectors corresponding to small eigenvalues in magnitude, e.g., the GMRES-E method \BookCite{134}. The basic ideas of the LGMRES method and the GMRES-E method are described next.}
{The group of $m$ iteration steps (e.g., lines 3 and 4 in Algorithm\ \ref{alg:gmres-restarted}) is called a cycle, and the approximate solution of the LGMRES method after the $j$th cycle is denoted by $\bm x_m^{(j)}$. Let $\bm z^{(k)}:=\bm x_m^{(j-1)}-\bm x_m^{(j-2)}$. Then for a given $l$, the LGMRES method finds an approximate solution in the following way:}
\[\bm x_m^{(j)}=\bm x_m^{(j-1)}+q_{m-1}^{(j)}(A)\bm r_m^{(j-1)}+\sum_{k=j-l}^{j-1}\alpha_{j,k}(A)\bm z^{(k)},\]
{where $q_{m-1}^{(j)}(A)$ is a polynomial of degree $m-1$, and $\alpha_{j,k}(A)$'s are polynomials of degree $k$. These polynomials are determined by minimizing the residual 2-norm of $\bm x_m^{(j)}$. For the related algorithms, see, e.g., \BookCite{13}, \BookCite{103}, \BookCite{104}, \BookCite{106}.}
% block: gmres-e-augmentation page:127 type:ordinary-paragraph
{The approximate solution of the GMRES-E method after the $j$th cycle is denoted by $\bm x_m^{(j)}$, and $\bm p_k^{(j)}$ is an approximate eigenvector (Harmonic Ritz vectors) that corresponds to the $k$th smallest approximate eigenvalue (Harmonic Ritz value) in magnitude. Then, for a given $p$ the GMRES-E method finds an approximate solution in the following way:}
\[\bm x_m^{(j)}=\bm x_m^{(j-1)}+q_{m-1}^{(j)}(A)\bm r_m^{(j-1)}+\sum_{k=1}^{p}\beta_{j,k}(A)\bm p_k^{(j-1)},\]
{where $q_{m-1}^{(j)}(A)$ is a polynomial of degree $m-1$, and $\beta_{j,k}(A)$'s are polynomials of degree $k$. These polynomials are determined by minimizing the residual 2-norm of $\bm x_m^{(j)}$. For the related algorithms, see, e.g., \BookCite{11}, \BookCite{135}. By combining the LGMRES method and the GMRES-E method, the LGMRES-E method using a switching controller is proposed in \BookCite{31}. Roughly speaking, the switching rule is to choose the LGMRES method while good convergence appears and to choose the GMRES-E method when slow convergence is found.}
{The augmentation is not only for the GMRES method. In fact, frameworks of augmentation and deflation for Krylov subspace methods including the BiCG method, the BiCR method, and the other important Krylov space methods are developed in \BookCite{74}, \BookCite{89}, \BookCite{90}.}
% block: gmres-numerical-example pages:127,128 type:ordinary-paragraph
{Finally, we provide the result of a numerical example for the GMRES method and the GMRES($m$) method. Consider nonsymmetric linear systems \eqref{eq:three-dimensional-difference-system} with parameters $N=20$, $a_1=a_2=a_3=1$, $b_1=b_2=b_3=1$, and $c=1$. The result is shown in Fig.\ \ref{fig:gmres-restarts}, where the stopping criterion used for the GMRES method and the GMRES($m$) method is $\|\bm r_n\|/\|\bm b\|\le10^{-10}$, and the horizontal axis is the number of matrix--vector multiplications.}
{We see from Fig.\ \ref{fig:gmres-restarts} that the residual 2-norms of the GMRES method, the GMRES(10) method, and the GMRES(20) method decrease monotonically, which are theoretically guaranteed. In terms of the number of matrix--vector multiplications, the GMRES method performs better than the restarted GMRES method (GMRES(10), GMRES(20)). As for the GMRES method, the required number of matrix--vector multiplications is the smallest of all, since it finds the best approximate solution over the Krylov subspace. On the other hand, as mentioned before, the computational cost of the GMRES method grows linearly as the number of iterations increases. This means that matrix--vector multiplications may not become the most time-consuming parts when the required number of iterations is very large. In terms of CPU time, the GMRES method required 0.39 sec., GMRES(10) 0.31 sec., GMRES(20) 0.26 sec.}\footnote{{Codes: GNU Octave version 6.2.0, OS: macOS Moterey version 12.2.1, CPU: Apple M1 pro.}}
{Therefore, with respect to CPU time, the GMRES(20) method was the best of all, which balanced the increasing computational costs and the required number of iterations.}
% block: gmres-restart-figure page:128 type:figure
\begin{figure}[H]\centering\includegraphics[width=.87\linewidth]{figures/gmres-restarts.png}
\caption{{Convergence histories of GMRES, GMRES(10), and GMRES(20)}}\label{fig:gmres-restarts}\end{figure}
% block: gcr-heading pages:128,129 type:subsection
\subsubsection{{The Generalized Conjugate Residual (GCR) Method}}\label{subsec:gcr-method}
{The Generalized Conjugate Residual (GCR) method \BookCite{54} is another Krylov solver based on the Arnoldi process. In exact precision arithmetic, the GCR method produces the same approximate solution as the GMRES method. The algorithm is obtained by the following Arnoldi process with the weight $A^{\mathrm H}A$ and $\bm r_0=\bm b-A\bm x_0$:}
{set} $\bm v_1=\dfrac{\bm r_0}{\|\bm r_0\|_{A^{\mathrm H}A}}$,\par
{for $n=1,2,\ldots$, do:}
\[h_{i,n}=(\bm v_i,A\bm v_n)_{A^{\mathrm H}A},\quad i=1,2,\ldots,n,\]
\begin{equation}\label{eq:gcr-weighted-arnoldi-remainder}\widetilde{\bm v}_{n+1}=A\bm v_n-\sum_{i=1}^nh_{i,n}\bm v_i,\end{equation}
\[h_{n+1,n}=\|\widetilde{\bm v}_{n+1}\|_{A^{\mathrm H}A},\]
\begin{equation}\label{eq:gcr-weighted-arnoldi-normalization}\bm v_{n+1}=\frac{\widetilde{\bm v}_{n+1}}{h_{n+1,n}}.\end{equation}
{end}
{The process generates $A^{\mathrm H}A$-orthonormal vectors, i.e.,}
\[(\bm v_i,A^{\mathrm H}A\bm v_j)=(A\bm v_i,A\bm v_j)=\delta_{i,j}.\]
% block: gcr-residual-expansion page:129 type:ordinary-paragraph
{Let $V_n$ be a matrix whose columns are $\bm v_1,\bm v_2,\ldots,\bm v_n$. Then, the GCR method generates}
\[\bm x_n=\bm x_0+V_n\bm y_n\]
{such that the residual 2-norm $\|\bm r_n\|$ is minimized. The corresponding residual vector can be written as}
\begin{equation}\label{eq:gcr-residual-expansion}\bm r_n=\bm r_0-AV_n\bm y_n=\bm r_0-\sum_{i=1}^ny_iA\bm v_i.\end{equation}
{From the above relation, we have}
\begin{equation}\label{eq:gcr-residual-basis-update}\bm r_n=\bm r_{n-1}-y_nA\bm v_n.\end{equation}
{Unknown parameters $y_i$ ($i=1,2,\ldots,n$) are determined by minimizing the residual 2-norm, i.e., from \eqref{eq:gcr-residual-expansion} this can be achieved by}
\[y_i=(A\bm v_i,\bm r_n),\quad\text{for }i=1,2,\ldots,n.\]
{Since $A\bm v_i$ is an orthonormal vector, we have}
\[y_i=(A\bm v_i,\bm r_0),\quad\text{for }i=1,2,\ldots,n.\]
{Here, we define the vector $\bm p_n:=-y_n\widetilde{\bm v}_{n+1}$. Then, it follows from \eqref{eq:gcr-weighted-arnoldi-remainder} that}
% block: gcr-direction-projection page:130 type:ordinary-paragraph
\begin{equation}\label{eq:gcr-direction-from-arnoldi}\bm p_n=-y_nA\bm v_n-\sum_{i=1}^n(\bm v_i,-y_nA\bm v_n)_{A^{\mathrm H}A}\bm v_i.\end{equation}
{From \eqref{eq:gcr-residual-basis-update} we have $-y_nA\bm v_n=\bm r_n-\bm r_{n-1}$. Substituting this relation into \eqref{eq:gcr-direction-from-arnoldi} yields}
\begin{equation}\label{eq:gcr-direction-residual-difference}\bm p_n=\bm r_n-\bm r_{n-1}-\sum_{i=1}^n(\bm v_i,\bm r_n-\bm r_{n-1})_{A^{\mathrm H}A}\bm v_i\end{equation}
\[=\bm r_n-\sum_{i=1}^n(\bm v_i,\bm r_n)_{A^{\mathrm H}A}\bm v_i-\left(\bm r_{n-1}-\sum_{i=1}^n(\bm v_i,\bm r_{n-1})_{A^{\mathrm H}A}\bm v_i\right).\]
{Since $\bm r_{n-1}$ lies in $\mathcal K_n(A,\bm r_0)$, it can be expanded as $\bm r_{n-1}=\sum_{i=1}^nc_i\bm v_i$. Then,}
\[(\bm v_j,\bm r_{n-1})_{A^{\mathrm H}A}=\sum_{i=1}^nc_i(\bm v_j,\bm v_i)=c_i\]
{because $(\bm v_j,\bm v_i)=\delta_{j,i}$. Thus we have}
\begin{equation}\label{eq:gcr-previous-residual-projection}\bm r_{n-1}=\sum_{i=1}^n(\bm v_i,\bm r_{n-1})_{A^{\mathrm H}A}\bm v_i.\end{equation}
{From \eqref{eq:gcr-direction-residual-difference} and \eqref{eq:gcr-previous-residual-projection}, the vector $\bm p_n$ can be written as}
\begin{equation}\label{eq:gcr-direction-weighted-projection}\bm p_n=\bm r_n-\sum_{i=1}^n(\bm v_i,\bm r_n)_{A^{\mathrm H}A}\bm v_i.\end{equation}
% block: gcr-normalized-direction page:130 type:ordinary-paragraph
{On the other hand, $\bm p_n$ and $\bm v_{n+1}$ are related by}
\[\bm p_n=-\frac{y_n}{|y_n|}\|A\bm p_n\|\bm v_{n+1},\]
{because it follows from the definition of $\bm p_n$ and \eqref{eq:gcr-weighted-arnoldi-normalization} that}
\[\begin{aligned}\bm p_n&=-y_n\widetilde{\bm v}_{n+1}\\&=-y_n\|\widetilde{\bm v}_{n+1}\|_{A^{\mathrm H}A}\bm v_{n+1}\\&=-y_n\sqrt{(A\widetilde{\bm v}_{n+1},A\widetilde{\bm v}_{n+1})}\,\bm v_{n+1}\\&=-y_n\sqrt{\left(-\frac1{y_n}A\bm p_n,-\frac1{y_n}A\bm p_n\right)}\,\bm v_{n+1}\\&=-\frac{y_n}{|y_n|}\sqrt{(A\bm p_n,A\bm p_n)}\,\bm v_{n+1}.\end{aligned}\]
% block: gcr-coefficients page:131 type:ordinary-paragraph
{Hence, using $\bm v_n=-(|y_{n-1}|/y_{n-1})\times\bm p_{n-1}/\sqrt{(A\bm p_{n-1},A\bm p_{n-1})}$, we have the following relations:}
\[\begin{aligned}y_nA\bm v_n&=(A\bm v_n,\bm r_0)A\bm v_n=\frac{(A\bm p_{n-1},\bm r_0)}{(A\bm p_{n-1},A\bm p_{n-1})}A\bm p_{n-1},\\(\bm v_i,\bm r_n)_{A^{\mathrm H}A}\bm v_i&=\frac{(A\bm p_{i-1},A\bm r_n)}{(A\bm p_{i-1},A\bm p_{i-1})}\bm p_{i-1}.\end{aligned}\]
{Therefore, substituting the above results into \eqref{eq:gcr-residual-basis-update} and \eqref{eq:gcr-direction-weighted-projection}, we have}
\[\begin{aligned}\bm r_n&=\bm r_{n-1}-\frac{(A\bm p_{n-1},\bm r_0)}{(A\bm p_{n-1},A\bm p_{n-1})}A\bm p_{n-1},\\\bm p_n&=\bm r_n-\sum_{i=1}^n\frac{(A\bm p_{i-1},A\bm r_n)}{(A\bm p_{i-1},A\bm p_{i-1})}\bm p_{i-1}.\end{aligned}\]
{Now, from the definition of $\bm p_n$, we have $\bm p_{n-1}=c\bm v_n$ for a constant $c$. Thus $\bm p_{n-1}$ is $A^{\mathrm H}A$-orthogonal to $\bm v_1,\ldots,\bm v_{n-1}$. From \eqref{eq:gcr-residual-expansion}, it follows that}
\[(A\bm p_{n-1},\bm r_{n-1})=(A\bm p_{n-1},\bm r_0)-\sum_{i=1}^{n-1}y_i(A\bm p_{n-1},A\bm v_i)=(A\bm p_{n-1},\bm r_0).\]
{Now, let us define $\alpha_{n-1}$ and $\beta_{n-1,i}$ as}
\begin{equation}\label{eq:gcr-alpha}\alpha_{n-1}=\frac{(A\bm p_{n-1},\bm r_{n-1})}{(A\bm p_{n-1},A\bm p_{n-1})},\end{equation}
\begin{equation}\label{eq:gcr-beta}\beta_{n-1,i}=-\frac{(A\bm p_{i-1},A\bm r_n)}{(A\bm p_{i-1},A\bm p_{i-1})}.\end{equation}
{Then, we have}
\begin{equation}\label{eq:gcr-residual-recurrence}\bm r_n=\bm r_{n-1}-\alpha_{n-1}A\bm p_{n-1},\end{equation}
\begin{equation}\label{eq:gcr-direction-recurrence}\bm p_n=\bm r_n+\sum_{i=1}^n\beta_{n-1,i}\bm p_{i-1}.\end{equation}
{From the recurrence relation of the residual vector, we obtain}
\begin{equation}\label{eq:gcr-solution-recurrence}\bm x_n=\bm x_{n-1}+\alpha_{n-1}\bm p_{n-1}.\end{equation}
{To reduce the number of matrix--vector multiplications, use the following recurrences:}
% block: gcr-matvec-recurrence page:132 type:ordinary-paragraph
\begin{equation}\label{eq:gcr-matvec-recurrence}A\bm p_n=A\bm r_n+\sum_{i=1}^n\beta_{n-1,i}A\bm p_{i-1}.\end{equation}
{From \eqref{eq:gcr-alpha}--\eqref{eq:gcr-matvec-recurrence}, the algorithm of the GCR method is obtained in Algorithm\ \ref{alg:gcr}.}
% block: gcr-algorithm page:132 type:ordinary-paragraph
\begin{algorithm}[H]\caption{{The GCR method}}\label{alg:gcr}\begin{algorithmic}[1]
\Statex\textbf{Input:} $\bm x_0\in\mathbb C^N,\ \beta_{-1}=0,\ \bm r_0=\bm b-A\bm x_0,\ \bm p_0=\bm r_0$
\Statex\textbf{Output:} $\bm x_n$
\For{$n=0,1,2,\ldots$, {until convergence}}
\State $\alpha_n=\frac{(A\bm p_n,\bm r_n)}{(A\bm p_n,A\bm p_n)}$
\State $\bm x_{n+1}=\bm x_n+\alpha_n\bm p_n$
\State $\bm r_{n+1}=\bm r_n-\alpha_nA\bm p_n$
\State $\beta_{n,i}=-\frac{(A\bm p_{i-1},A\bm r_{n+1})}{(A\bm p_{i-1},A\bm p_{i-1})},\quad1\le i\le n+1$
\State $\bm p_{n+1}=\bm r_{n+1}+\sum_{i=1}^{n+1}\beta_{n,i}\bm p_{i-1}$
\State $(A\bm p_{n+1}=A\bm r_{n+1}+\sum_{i=1}^{n+1}\beta_{n,i}A\bm p_{i-1})$
\EndFor\end{algorithmic}\end{algorithm}
% block: gcr-properties-proposition page:132 type:proposition
{The properties of the GCR method are given next.}
\par\Needspace{3\baselineskip}
\begin{Proposition}\label{prop:gcr-orthogonality}
{The GCR iterates hold the following properties:}
\begin{enumerate}[label=(G\arabic*),ref=G\arabic*]
\item\label{part:gcr-direction-orthogonality} $(A\bm p_i,A\bm p_j)=0,\quad i\ne j$,
\item $(\bm r_i,A\bm p_j)=0,\quad i>j$,
\item $(\bm r_i,A\bm p_i)=(\bm r_i,A\bm r_i)$,
\item $(\bm r_i,A\bm r_j)=0,\quad i>j$,
\item $(\bm r_i,A\bm p_i)=(\bm r_0,A\bm p_i),\quad i\ge j$.
\end{enumerate}\end{Proposition}
{For the proof, see \BookCite{54}, Theorem 3.1.}
% block: gcr-restarted-and-orthomin page:132 type:ordinary-paragraph
{From the above properties, we see that the previous derivation is based on the property (\ref{part:gcr-direction-orthogonality}). In exact precision arithmetic, the GCR method, as well as the GMRES method, converges within at most $N$ iterations. However, the GCR method has the same shortcoming as the GMRES method in that the computational work and the required memory increase with the number of iterations. Hence, a restarted version of GCR which is referred to as the GCR($k$) method was proposed in \BookCite{54}. The algorithm is described in Algorithm\ \ref{alg:gcr-restarted}.}
{Another alternative is to keep only $k$-direction vectors:}
\[\bm p_{n+1}=\bm r_{n+1}+\sum_{i=1}^{\min\{k,n+1\}}\beta_{n,i}\bm p_{i-1}.\]
{This method is known as the Orthomin method proposed by Vinsome \BookCite{201}. Later the Orthomin method was referred to as the Orthomin($k$) method in \BookCite{54}. The algorithm of the Orthomin($k$) method is described in Algorithm\ \ref{alg:orthomin}.}
% block: gcr-restarted-algorithm page:133 type:ordinary-paragraph
\begin{algorithm}[H]\caption{{The GCR($k$) method}}\label{alg:gcr-restarted}\begin{algorithmic}[1]
\Statex\textbf{Input:} $\bm x_0\in\mathbb C^N,\ \bm r_0=\bm b-A\bm x_0,\ \bm p_0=\bm r_0$
\Statex\textbf{Output:} $\bm x_{k+1}$
\For{$n=0,1,\ldots,k$ {until convergence}}
\State $\alpha_n=\frac{(A\bm p_n,\bm r_n)}{(A\bm p_n,A\bm p_n)}$
\State $\bm x_{n+1}=\bm x_n+\alpha_n\bm p_n$
\State $\bm r_{n+1}=\bm r_n-\alpha_nA\bm p_n$
\State $\beta_{n,i}=-\frac{(A\bm p_{i-1},A\bm r_{n+1})}{(A\bm p_{i-1},A\bm p_{i-1})},\quad1\le i\le n+1$
\State $\bm p_{n+1}=\bm r_{n+1}+\sum_{i=1}^{n+1}\beta_{n,i}\bm p_{i-1}$
\State $(A\bm p_{n+1}=A\bm r_{n+1}+\sum_{i=1}^{n+1}\beta_{n,i}A\bm p_{i-1})$
\EndFor\State $\bm x_0=\bm x_{k+1}$\State{repeat}
\end{algorithmic}\end{algorithm}
% block: orthomin-algorithm page:133 type:ordinary-paragraph
\begin{algorithm}[H]\caption{{The Orthomin($k$) method}}\label{alg:orthomin}\begin{algorithmic}[1]
\Statex\textbf{Input:} $\bm x_0\in\mathbb C^N,\ \bm r_0=\bm b-A\bm x_0,\ \bm p_0=\bm r_0$
\Statex\textbf{Output:} $\bm x_n$
\For{$n=0,1,\ldots$, {until $\|\bm r_n\|\le\epsilon\|\bm b\|$}}
\State $\alpha_n=\frac{(A\bm p_n,\bm r_n)}{(A\bm p_n,A\bm p_n)}$
\State $\bm x_{n+1}=\bm x_n+\alpha_n\bm p_n$
\State $\bm r_{n+1}=\bm r_n-\alpha_nA\bm p_n$
\State $\beta_{n,i}=-\frac{(A\bm p_{i-1},A\bm r_{n+1})}{(A\bm p_{i-1},A\bm p_{i-1})},\quad n-k+2\le i\le n+1$
\State $\bm p_{n+1}=\bm r_{n+1}+\sum_{i=1}^{\min\{k,n+1\}}\beta_{n,i}\bm p_{i-1}$
\State $(A\bm p_{n+1}=A\bm r_{n+1}+\sum_{i=1}^{n+1}\beta_{n,i}A\bm p_{i-1})$
\EndFor\end{algorithmic}\end{algorithm}
% block: gcr-convergence-proposition pages:133,134 type:proposition
{The convergence analysis of the GCR method, the GCR($k$) method, and the Orthomin($k$) method is given in Proposition\ \ref{prop:gcr-restarted-convergence}.}
\par\Needspace{3\baselineskip}
\begin{Proposition}\label{prop:gcr-restarted-convergence}
{Let $M$ be the symmetric part of $A\in\mathbb R^{N\times N}$ and let $R$ be the skew-symmetric part of $A$, i.e., $M=(A+A^{\top})/2$, $R=(A-A)^{\top}/2$. If $M$ is positive definite, then $\{\bm r_n\}$ that are the sequence of residuals generated by the GCR method, the GCR($k$) method, or the Orthomin($k$) method satisfy}
\[\|\bm r_n\|\le\left(1-\frac{\lambda_{\min}(M)^2}{\lambda_{\max}(A^{\top}A)}\right)^{n/2}\|\bm r_0\|\]
{and}
\[\|\bm r_n\|\le\left(1-\frac{\lambda_{\min}(M)^2}{\lambda_{\min}(M)\lambda_{\max}(M)+\rho(R)^2}\right)^{n/2}\|\bm r_0\|,\]
{where $\lambda_{\min}(M)$ and $\lambda_{\max}(M)$ are the minimum and maximum eigenvalues of the symmetric positive definite matrix $M$, and $\rho(R)$ is the spectral radius of $R$.}
\end{Proposition}
{See \eqref{eq:stationary-spectral-radius} for the spectral radius. For the proof, see \BookCite{54}, Theorem 4.4.}
{We can see from Proposition\ \ref{prop:gcr-restarted-convergence} that if $A$ is close to $I$, then the three methods converge fast since $\lambda_{\min}(M)\approx\lambda_{\max}(M)\approx1$ and $R\approx O$.}
% block: fom-heading page:134 type:subsection
\subsubsection{{The Full Orthogonalization Method (FOM)}}\label{subsec:fom}
{The FOM \BookCite{149} is a generalization of the CG method. As well as the CG method, the FOM finds approximate solutions such that}
\[\begin{aligned}\bm x_n&=\bm x_0+\bm z_n,\quad\bm z_n\in\mathcal K_n(A,\bm r_0),\\\bm r_n&\perp\mathcal K_n(A,\bm r_0).\end{aligned}\]
{The FOM uses the Arnoldi method in Sect.\ \ref{subsec:arnoldi-process} to produce the orthonormal basis vectors of $\mathcal K_n(A,\bm r_0)$. Thus the residual vector $\bm r_n$ has the same form as \eqref{eq:gmres-projected-residual}, i.e.,}
\begin{equation}\label{eq:fom-projected-residual}\bm r_n=V_{n+1}(\beta\bm e_1-H_{n+1,n}\bm y_n).\end{equation}
{The difference between the GMRES method and the FOM is that the GMRES method finds $\bm y_n$ such that the 2-norm of the residual vector $\|\bm r_n\|$ is minimized, and the FOM finds $\bm y_n$ such that the residual vector is orthogonal to the $n$ dimensional Krylov subspace $\bm r_n\perp\mathcal K_n(A,\bm r_0)$. Let $\bm v_i$ be the $i$th column of $V_{n+1}$. Then}
\[\begin{aligned}V_n^{\mathrm H}\bm r_n=\bm0&\Leftrightarrow\bm v_i^{\mathrm H}\bm r_n=0\quad\text{for }i=1,2,\ldots,n\\&\Leftrightarrow\bm r_n^{\mathrm H}(c_1\bm v_1+c_2\bm v_2+\cdots+c_n\bm v_n)=0\quad\text{for all }c_i\in\mathbb C\\&\Leftrightarrow\bm r_n\perp\mathcal K_n(A,\bm r_0).\end{aligned}\]
% block: fom-projected-solve page:134 type:ordinary-paragraph
{Note that the statement $\bm v_i^{\mathrm H}\bm r_n=0\Leftarrow\bm r_n^{\mathrm H}(c_1\bm v_1+c_2\bm v_2+\cdots+c_n\bm v_n)=0$ can be shown by setting $c_i=1$ and $c_j=0$ for $i\ne j$. From this, finding $\bm y_n$ such that $\bm r_n\perp\mathcal K_n(A,\bm r_0)$ is equivalent to determining $\bm y_n$ such that $V_n^{\mathrm H}\bm r_n=\bm0$. Then, from \eqref{eq:fom-projected-residual} it follows that}
\[\bm0=V_n^{\mathrm H}\bm r_n=V_n^{\mathrm H}V_{n+1}(\beta\bm e_1-H_{n+1,n}\bm y_n)=\beta\bm e_1-H_{n,n}\bm y_n.\]
{Thus $\bm y_n$ can be obtained by solving the following linear systems:}
\[H_{n,n}\bm y_n=\beta\bm e_1.\]
{Since $H_{n,n}$ is a Hessenberg matrix, the linear systems can be solved by the LU decomposition with the computational cost of $O(n^2)$. After the advent of the GMRES method, the FOM became less attractive because, unlike the GMRES method, the FOM may suffer from breakdown due to using the LU decomposition without pivoting. On the other hand, the FOM is still attractive to use for solving shifted linear systems in Chap.\ \ref{chap:shifted-linear-systems}, when considering the restart.}
% block: product-type-heading page:135 type:subsection
\subsubsection{{Product-Type Krylov Subspace Methods}}\label{subsec:product-type-krylov}
{As seen in Sect.\ \ref{subsec:bicg-method}, the BiCG method requires $A^{\mathrm H}$ for solving $A\bm x=\bm b$. P. Sonneveld found that squaring the BiCG polynomials leads to an algorithm without using the information of $A^{\mathrm H}$, and the resulting algorithm is known as the CGS method \BookCite{173}. Though the convergence depends on linear systems, the CGS method is much faster than the BiCG method when the BiCG method shows smooth convergence behavior. However, the CGS method sometimes has much irregular convergence and may suffer from loss of accuracy of the approximate solution or may diverge. H. A. van der Vorst observed that squaring the BiCG polynomials can be replaced by the multiplication of degree-one minimum residual polynomials and the BiCG polynomial, leading to the BiCGSTAB method \BookCite{195}, which tends to have smoother convergence behavior than the CGS method. M. H. Gutknecht proposed the multiplication of degree-two minimum residual polynomials and the BiCG polynomials, which is known as the BiCGSTAB2 method \BookCite{87}. G. L. G. Sleijpen and D. R. Fokkema proposed the multiplication of degree-$\ell$ minimal residual polynomials and the BiCG polynomials, which is known as the BiCGSTAB($\ell$) method \BookCite{162}. S.-L. Zhang constructed a framework that includes the CGS method, the BiCGSTAB method, and the BiCGSTAB2 method. In the framework, the GPBiCG method was proposed in \BookCite{212}. These methods are referred to as product-type methods (based on the Bi-CG method) first named by S.-L. Zhang.}
{In this section, Zhang's framework is explained, based on Zhang's book in \BookCite{72}.}
% block: product-bicg-polynomials page:135 type:ordinary-paragraph
{As seen in \eqref{eq:cg-residual-direction-polynomials}, the residual vector of the CG method can be written as the multiplication of Lanczos polynomials and the initial residual vector $\bm r_0$. Similarly, the residual vector of the BiCG method can also be written as follows:}
\[\bm r_n^{\mathrm{BiCG}}=R_n(A)\bm r_0,\]
{where}
\begin{equation}\label{eq:product-bicg-polynomials-initial}R_0(\lambda)=1,\quad P_0(\lambda)=1,\end{equation}
\begin{equation}\label{eq:product-bicg-residual-polynomial}R_n(\lambda)=R_{n-1}(\lambda)-\alpha_{n-1}\lambda P_{n-1}(\lambda),\end{equation}
\begin{equation}\label{eq:product-bicg-direction-polynomial}P_n(\lambda)=R_n(\lambda)+\beta_{n-1}P_{n-1}(\lambda)\quad n=1,2,\ldots\end{equation}
{The differences between the (matrix) polynomial representation of the CG method and that of the BiCG method are $\alpha_i$ and $\beta_i$ of the Lanczos polynomials. The $n$th residual vector of product-type methods based on the BiCG method is defined by the product of $n$ degree polynomials and the $n$th residual vector of the BiCG method as follows:}
\begin{equation}\label{eq:product-residual-polynomial}\bm r_n=H_n(A)\bm r_n^{\mathrm{BiCG}}=H_n(A)R_n(A)\bm r_0.\end{equation}
% block: product-polynomial-restructuring page:136 type:subsection
\paragraph{{Restructuring of Residual Polynomials}}\label{subsec:product-polynomial-restructuring}\mbox{}\par
{S.-L. Zhang proposed $H_n$ in \eqref{eq:product-residual-polynomial} as the following three-term recurrence relations:}
\begin{equation}\label{eq:product-h-initial}H_0(\lambda):=1,\end{equation}
\begin{equation}\label{eq:product-h-first}H_1(\lambda):=(1-\zeta_0\lambda)H_0(\lambda),\end{equation}
\begin{equation}\label{eq:product-h-three-term}H_n(\lambda):=(1+\eta_{n-1}-\zeta_{n-1}\lambda)H_{n-1}(\lambda)-\eta_{n-1}H_{n-2}(\lambda),\quad n=2,3,\ldots\end{equation}
{We see that $H_n$ is similar to the Lanczos polynomials \eqref{eq:cg-polynomial-initial-zero}--\eqref{eq:cg-polynomial-three-term}.}
{In the following, $H_n$ is rewritten as the form of coupled two-term recurrence relation. We introduce $G_{n-1}$ as}
\[G_{n-1}(\lambda):=\frac{H_{n-1}(\lambda)-H_n(\lambda)}{\lambda},\]
{and we rewrite \eqref{eq:product-h-three-term} as}
\[-(H_{n-1}(\lambda)-H_n(\lambda))=-\zeta_{n-1}\lambda H_{n-1}(\lambda)-\eta_{n-1}(H_{n-2}(\lambda)-H_{n-1}(\lambda)).\]
{Then $H_n(\lambda)$ in \eqref{eq:product-h-initial}--\eqref{eq:product-h-three-term} can be rewritten as}
\begin{equation}\label{eq:product-coupled-initial}H_0(\lambda)=1,\quad G_0(\lambda)=\zeta_0,\end{equation}
\begin{equation}\label{eq:product-coupled-h}H_n(\lambda)=H_{n-1}(\lambda)-\lambda G_{n-1}(\lambda),\end{equation}
\begin{equation}\label{eq:product-coupled-g}G_n(\lambda)=\zeta_nH_n(\lambda)+\eta_nG_{n-1}(\lambda),\quad n=1,2,\ldots\end{equation}
% block: product-iterate-polynomial-recurrences page:136 type:subsection
\paragraph{{Recurrence Formulas for the Iterates of the Product-Type Methods}}\label{subsec:product-iterate-recurrences}\mbox{}\par
{Since recurrence relations of $R_n,P_n,H_n,G_n$ are already given in \eqref{eq:product-bicg-polynomials-initial}--\eqref{eq:product-bicg-direction-polynomial} and in \eqref{eq:product-coupled-initial}--\eqref{eq:product-coupled-g}, the residual vector $\bm r_{n+1}=H_{n+1}(A)R_{n+1}(A)\bm r_0$ can be computed by the previous residual vector $\bm r_n=H_n(A)R_n(A)\bm r_0$ using the recurrence relations. In fact, $H_{n+1}R_{n+1}$ can be expanded as follows:}
\begin{equation}\label{eq:product-polynomial-residual-three}H_{n+1}R_{n+1}=H_nR_{n+1}-\eta_n\lambda G_{n-1}R_{n+1}-\zeta_n\lambda H_nR_{n+1}\end{equation}
\begin{equation}\label{eq:product-polynomial-residual-two}=H_nR_n-\alpha_n\lambda H_nP_n-\lambda G_nR_{n+1},\end{equation}
\begin{equation}\label{eq:product-polynomial-t}H_nR_{n+1}=H_nR_n-\alpha_n\lambda H_nP_n,\end{equation}
\begin{equation}\label{eq:product-polynomial-y}\lambda G_nR_{n+2}=H_nR_{n+1}-H_{n+1}R_{n+1}-\alpha_{n+1}\lambda H_nP_{n+1}+\alpha_{n+1}\lambda H_{n+1}P_{n+1},\end{equation}
\begin{equation}\label{eq:product-polynomial-p}H_{n+1}P_{n+1}=H_{n+1}R_{n+1}+\beta_nH_nP_n-\beta_n\lambda G_nP_n,\end{equation}
\begin{equation}\label{eq:product-polynomial-w}\lambda H_nP_{n+1}=\lambda H_nR_{n+1}+\beta_n\lambda H_nP_n,\end{equation}
\begin{equation}\label{eq:product-polynomial-u}\lambda G_nP_n=\zeta_n\lambda H_nP_n+\eta_n(H_{n-1}R_n-H_nR_n+\beta_{n-1}\lambda G_{n-1}P_{n-1}),\end{equation}
\begin{equation}\label{eq:product-polynomial-z}G_nR_{n+1}=\zeta_nH_nR_n+\eta_nG_{n-1}R_n-\alpha_n\lambda G_nP_n.\end{equation}
% block: product-auxiliary-vector-recurrences page:137 type:ordinary-paragraph
{Let us define new auxiliary vectors as}
\[\begin{aligned}\bm t_n&:=H_n(A)\bm r_{n+1}^{\mathrm{BiCG}},&\bm y_n&:=AG_{n-1}(A)\bm r_{n+1}^{\mathrm{BiCG}},\\\bm p_n&:=H_n(A)\bm p_n^{\mathrm{BiCG}},&\bm w_n&:=AH_n(A)\bm p_{n+1}^{\mathrm{BiCG}},\\\bm u_n&:=AG_n(A)\bm p_n^{\mathrm{BiCG}},&\bm z_n&:=G_n(A)\bm r_{n+1}^{\mathrm{BiCG}}.\end{aligned}\]
{Then, from \eqref{eq:product-polynomial-residual-three}--\eqref{eq:product-polynomial-z} we have the following recurrence formulas:}
\begin{equation}\label{eq:product-residual-three}\bm r_{n+1}=\bm t_n-\eta_n\bm y_n-\zeta_nA\bm t_n\end{equation}
\begin{equation}\label{eq:product-residual-two}=\bm r_n-\alpha_nA\bm p_n-A\bm z_n,\end{equation}
\begin{equation}\label{eq:product-t}\bm t_n=\bm r_n-\alpha_nA\bm p_n,\end{equation}
\begin{equation}\label{eq:product-y}\bm y_{n+1}=\bm t_n-\bm r_{n+1}-\alpha_{n+1}\bm w_n+\alpha_{n+1}A\bm p_{n+1},\end{equation}
\begin{equation}\label{eq:product-p}\bm p_{n+1}=\bm r_{n+1}+\beta_n(\bm p_n-\bm u_n),\end{equation}
\begin{equation}\label{eq:product-w}\bm w_n=A\bm t_n+\beta_nA\bm p_n,\end{equation}
\begin{equation}\label{eq:product-u}\bm u_n=\zeta_nA\bm p_n+\eta_n(\bm t_{n-1}-\bm r_n+\beta_{n-1}\bm u_{n-1}),\end{equation}
\begin{equation}\label{eq:product-z}\bm z_n=\zeta_n\bm r_n+\eta_n\bm z_{n-1}-\alpha_n\bm u_n.\end{equation}
{Two recurrence relations, \eqref{eq:product-residual-three} and \eqref{eq:product-residual-two}, are described for the residual vector and the approximate solution. The first recurrence relation \eqref{eq:product-residual-three} will be used to determine the two parameters $\zeta_n$ and $\eta_n$, and the second recurrence relation \eqref{eq:product-residual-two} can be used to obtain the approximate solution as follows:}
\[\bm x_{n+1}=\bm x_n+\alpha_n\bm p_n+\bm z_n.\]
% block: product-alpha-beta page:137 type:subsection
\paragraph{{Computations for $\alpha_n$ and $\beta_n$}}\label{subsec:product-alpha-beta}\mbox{}\par
{Since the coefficient of the highest-order term of $H_n$ is $(-1)^n\prod_{i=0}^{n-1}\zeta_i$, we obtain}
\[\begin{aligned}(\bm r_0^{*},\bm r_n)&=(\bar H_n(A^{\mathrm H})\bm r_0^{*},\bm r_n^{\mathrm{BiCG}})=\left((-1)^n\prod_{i=0}^{n-1}\zeta_i\right)((A^{\mathrm H})^n\bm r_0^{*},\bm r_n^{\mathrm{BiCG}}),\\(\bm r_0^{*},A\bm p_n)&=(\bar H_n(A^{\mathrm H})\bm r_0^{*},A\bm p_n^{\mathrm{BiCG}})=\left((-1)^n\prod_{i=0}^{n-1}\zeta_i\right)((A^{\mathrm H})^n\bm r_0^{*},A\bm p_n^{\mathrm{BiCG}}).\end{aligned}\]
{Thus, from \eqref{eq:bicg-alpha-shadow-power} and \eqref{eq:bicg-beta-shadow-power} it follows that}
\[\alpha_n=\frac{(\bm r_0^{*},\bm r_n)}{(\bm r_0^{*},A\bm p_n)},\qquad\beta_n=\frac{\alpha_n}{\zeta_n}\times\frac{(\bm r_0^{*},\bm r_{n+1})}{(\bm r_0^{*},\bm r_n)}.\]
% block: product-implementation page:138 type:subsection
\paragraph{{Implementation Details}}\label{subsec:product-implementation}\mbox{}\par
{Krylov subspace methods, such as the CGS method, the BiCGSTAB method, the BiCGSTAB2 method, and the GPBiCG method, can be derived from Zhang's framework by determining the two parameters $\zeta_n$ and $\eta_n$. The relations among these Krylov subspace methods and the parameters are listed in Table\ \ref{tab:product-parameters}.}
% block: product-parameters-table page:138 type:table
\begin{table}[H]\centering
\caption{{Choices of $\zeta_n$ and $\eta_n$ for the product-type methods}}\label{tab:product-parameters}
\begin{tabular}{l|l}\hline
CGS&$\zeta_n=\alpha_n,\quad\eta_n=\dfrac{\beta_{n-1}}{\alpha_{n-1}}\alpha_n$\\\hline
BiCGSTAB&$\zeta_n=\operatorname*{arg\,min}_{\zeta_n\in\mathbb C}\|\bm r_{n+1}\|,\quad\eta_n=0$\\\hline
GPBiCG&$\{\zeta_n,\eta_n\}=\operatorname*{arg\,min}_{\zeta_n,\eta_n\in\mathbb C}\|\bm r_{n+1}\|$\\\hline
BiCGSTAB2&{BiCGSTAB parameters at even iterations}\\
&{GPBiCG parameters at odd iterations}\\\hline
\end{tabular}\end{table}
% block: cgs-parameter-choice page:138 type:subsection
\paragraph{{The Choice for the CGS Method}}\label{subsec:cgs-parameter-choice}\mbox{}\par
{Setting $\zeta_n=\alpha_n$ and $\eta_n=(\beta_{n-1}/\alpha_{n-1})\alpha_n$ in recurrence relations \eqref{eq:product-coupled-h} and \eqref{eq:product-coupled-g} yields the CGS method, i.e., the polynomials $H_n$ and $G_n$ become the Lanczos polynomials:}
\[H_n(\lambda)=R_n(\lambda),\qquad G_n(\lambda)=\alpha_nP_n(\lambda).\]
{This fact leads to the relation $\bm p_n-\bm u_n=\bm z_n/\alpha_n$, since from the definitions of auxiliary vectors we have}
\[\begin{aligned}\bm p_n-\bm u_n&=H_n(A)P_n(A)\bm r_0-AG_n(A)P_n(A)\bm r_0\\&=R_n(A)P_n(A)\bm r_0-\alpha_nAP_n(A)P_n(A)\bm r_0\\&=(R_n(A)-\alpha_nAP_n(A))P_n(A)\bm r_0\\&=R_{n+1}(A)P_n(A)\bm r_0\\&=P_n(A)R_{n+1}(A)\bm r_0\\&=P_n(A)\bm r_{n+1}^{\mathrm{BiCG}}\\&=\frac1{\alpha_n}\times G_n(A)\bm r_{n+1}^{\mathrm{BiCG}}\\&=\frac1{\alpha_n}\times\bm z_n.\end{aligned}\]
% block: cgs-auxiliary-rewrite page:139 type:ordinary-paragraph
{Notice that $\bm t_{n-1}-\bm r_n=A\bm z_{n-1}$ from \eqref{eq:product-residual-two} and \eqref{eq:product-t}. We use the recurrence formula \eqref{eq:product-residual-two} to update $\bm r_{n+1}$, and then the auxiliary vectors $\bm t_n,\bm y_n$, and $\bm w_n$ are not required in the recurrence formulas \eqref{eq:product-residual-three}--\eqref{eq:product-z}. Now we introduce the following auxiliary vectors:}
\[\widetilde{\bm u}_n:=A^{-1}\bm u_n/\alpha_n,\qquad\widetilde{\bm z}_n:=\bm z_n/\alpha_n,\]
{instead of $\bm u_n$ and $\bm z_n$. By using the relation $(\bm r_0^{*},A\bm p_n)=(\bm r_0^{*},\bm u_n/\alpha_n)$, the CGS method can be obtained and the algorithm is described in Algorithm\ \ref{alg:cgs}. Here $\widetilde{\bm u}_n$ and $\widetilde{\bm z}_n$ were rewritten as $\bm u_n$ and $\bm z_n$.}
% block: cgs-algorithm page:139 type:ordinary-paragraph
\begin{algorithm}[H]\caption{{The CGS method}}\label{alg:cgs}\begin{algorithmic}[1]
\Statex\textbf{Input:} $\bm x_0\in\mathbb C^N,\ \beta_{-1}=0,\ \bm u_{-1}=\bm z_{-1}=\bm0,\ \bm r_0=\bm b-A\bm x_0$
\Statex\textbf{Input:} {Choose $\bm r_0^{*}\in\mathbb C^N$, e.g., $\bm r_0^{*}=\bm r_0$}
\Statex\textbf{Output:} $\bm x_n$
\For{$n=0,1,\ldots$}
\State $\bm p_n=\bm r_n+\beta_{n-1}\bm z_{n-1}$
\State $\bm u_n=\bm p_n+\beta_{n-1}(\bm z_{n-1}+\beta_{n-1}\bm u_{n-1})$
\State $\alpha_n=\frac{(\bm r_0^{*},\bm r_n)}{(\bm r_0^{*},A\bm u_n)}$
\State $\bm z_n=\bm p_n-\alpha_nA\bm u_n$
\State $\bm x_{n+1}=\bm x_n+\alpha_n(\bm p_n+\bm z_n)$
\State $\bm r_{n+1}=\bm r_n-\alpha_nA(\bm p_n+\bm z_n)$
\State $\beta_n=\frac{(\bm r_0^{*},\bm r_{n+1})}{(\bm r_0^{*},\bm r_n)}$
\EndFor\end{algorithmic}\end{algorithm}
{The residual vector of the CGS method is written as $\bm r_{n+1}=R_{n+1}(A)R_{n+1}(A)\bm r_0$. The CGS method is short for the (Bi)CG squared method, which comes from the term $R_{n+1}(A)R_{n+1}(A)$. From this, the CGS method is expected to converge as twice as fast as the BiCG method. On the other hand, when the BiCG method shows irregular convergence behavior as is often the case, the CGS method gives much irregular convergence behavior, leading to much less accurate approximate solutions or no convergence.}
% block: bicgstab-parameter-choice pages:139,140 type:subsection
\paragraph{{The Choice for the BiCGSTAB Method}}\label{subsec:bicgstab-parameter-choice}\mbox{}\par
{If we choose $\eta_n=0$ and $\zeta_n$ such that the 2-norm of the residual vector $\|\bm r_{n+1}\|=\|\bm t_n-\zeta_nA\bm t_n\|$ is minimized, then the BiCGSTAB method is obtained. To achieve the minimization, we choose $\zeta_n$ such that the following property holds:}
\[\bm r_{n+1}\perp A\bm t_n,\]
{or equivalently it follows from \eqref{eq:scalar-residual-minimizer} that}
\[\zeta_n=\frac{\bm t_n^{\mathrm H}A\bm t_n}{(A\bm t_n)^{\mathrm H}(A\bm t_n)}.\]
{Since $\eta_n=0$, we obtain}
\[\bm u_n=\zeta_nA\bm p_n,\qquad\bm z_n=\zeta_n\bm t_n.\]
{The algorithm of the BiCGSTAB method is described in Algorithm\ \ref{alg:bicgstab}.}
{The BiCGSTAB method is short for the BiCG stabilized method. Below is an explanation of the meaning of the stabilization. Since $\eta_n=0$, the recurrence relations in \eqref{eq:product-coupled-h} and \eqref{eq:product-coupled-g} reduce to}
\[H_n(\lambda)=H_{n-1}(\lambda)-\zeta_{n-1}\lambda H_{n-1}(\lambda)=(1-\zeta_{n-1}\lambda)H_{n-1}(\lambda).\]
{This means that the residual vector of the BiCGSTAB method is given by}
\[\bm r_{n+1}=H_{n+1}(A)R_{n+1}(A)\bm r_0=(I-\zeta_nA)\bm r_n^{\mathrm{BiCG}}.\]
{This implies that since $\zeta_n$ is chosen such that $\|\bm r_{n+1}\|$ is minimized, the residual vector of the BiCG method is expected to be stabilized, i.e., the BiCGSTAB method is expected to have smoother convergence behavior than the BiCG method.}
% block: bicgstab-algorithm page:140 type:ordinary-paragraph
\begin{algorithm}[H]\caption{{The BiCGSTAB method}}\label{alg:bicgstab}\begin{algorithmic}[1]
\Statex\textbf{Input:} $\bm x_0\in\mathbb C^N,\ \beta_{-1}=0,\ \bm p_{-1}=\bm0,\ \bm r_0=\bm b-A\bm x_0$
\Statex\textbf{Input:} {Choose $\bm r_0^{*}\in\mathbb C^N$, e.g., $\bm r_0^{*}=\bm r_0$}
\Statex\textbf{Output:} $\bm x_n$
\For{$n=0,1,\ldots$}
\State $\bm p_n=\bm r_n+\beta_{n-1}(\bm p_{n-1}-\zeta_{n-1}A\bm p_{n-1})$
\State $\alpha_n=\frac{(\bm r_0^{*},\bm r_n)}{(\bm r_0^{*},A\bm p_n)}$
\State $\bm t_n=\bm r_n-\alpha_nA\bm p_n$
\State $\zeta_n=\frac{(A\bm t_n,\bm t_n)}{(A\bm t_n,A\bm t_n)}$
\State $\bm x_{n+1}=\bm x_n+\alpha_n\bm p_n+\zeta_n\bm t_n$
\State $\bm r_{n+1}=\bm t_n-\zeta_nA\bm t_n$
\State $\beta_n=\frac{\alpha_n}{\zeta_n}\times\frac{(\bm r_0^{*},\bm r_{n+1})}{(\bm r_0^{*},\bm r_n)}$
\EndFor\end{algorithmic}\end{algorithm}
% block: gpbicg-parameter-choice pages:140,141 type:subsection
\paragraph{{The Choice for the GPBiCG Method}}\label{subsec:gpbicg-parameter-choice}\mbox{}\par
{For the GPBiCG method, two parameters $\eta_n$ and $\zeta_n$ are chosen by minimizing the 2-norm of the residual vector $\|\bm r_{n+1}\|=\|\bm t_n-\eta_n\bm y_n-\zeta_nA\bm t_n\|$. The following orthogonal condition enables us to achieve the minimization:}
\[\bm r_{n+1}\perp A\bm t_n,\qquad\bm r_{n+1}\perp\bm y_n,\]
{or equivalently it follows from \eqref{eq:least-squares-decomposition} and \eqref{eq:least-squares-minimizer} with setting $\bm d=\bm t_n$, $M=[\bm y_n,A\bm t_n]$, $\bm x=[\eta_n,\zeta_n]^{\top}$ that we have}
\[\begin{aligned}\zeta_n&=\frac{(\bm y_n,\bm y_n)(A\bm t_n,\bm t_n)-(\bm y_n,\bm t_n)(A\bm t_n,\bm y_n)}{(A\bm t_n,A\bm t_n)(\bm y_n,\bm y_n)-(\bm y_n,A\bm t_n)(A\bm t_n,\bm y_n)},\\\eta_n&=\frac{(A\bm t_n,A\bm t_n)(\bm y_n,\bm t_n)-(\bm y_n,A\bm t_n)(A\bm t_n,\bm t_n)}{(A\bm t_n,A\bm t_n)(\bm y_n,\bm y_n)-(\bm y_n,A\bm t_n)(A\bm t_n,\bm y_n)}.\end{aligned}\]
{The algorithm of the GPBiCG method is described in Algorithm\ \ref{alg:gpbicg}. Note the GPBiCG method uses two parameters $\zeta_n$ and $\eta_n$ for the minimization of the residual 2-norm, and the BiCGSTAB method uses one parameter $\zeta_n$ for the minimization of the residual 2-norm and $\eta_n=0$. From this, the GPBiCG method is expected to give smoother convergence behavior than the BiCGSTAB method.}
% block: gpbicg-algorithm page:141 type:ordinary-paragraph
\begin{algorithm}[H]\caption{{The GPBiCG method}}\label{alg:gpbicg}\begin{algorithmic}[1]
\Statex\textbf{Input:} $\bm x_0\in\mathbb C^N,\ \beta_{-1}=0,\ \bm p_{-1}=\bm t_{-1}=\bm u_{-1}=\bm w_{-1}=\bm z_{-1}=\bm0,\ \bm r_0=\bm b-A\bm x_0$
\Statex\textbf{Input:} {Choose $\bm r_0^{*}\in\mathbb C^N$, e.g., $\bm r_0^{*}=\bm r_0$}
\Statex\textbf{Output:} $\bm x_n$
\For{$n=0,1,\ldots$}
\State $\bm p_n=\bm r_n+\beta_{n-1}(\bm p_{n-1}-\bm u_{n-1})$
\State $\alpha_n=\frac{(\bm r_0^{*},\bm r_n)}{(\bm r_0^{*},A\bm p_n)}$
\State $\bm y_n=\bm t_{n-1}-\bm r_n-\alpha_n\bm w_{n-1}+\alpha_nA\bm p_n$
\State $\bm t_n=\bm r_n-\alpha_nA\bm p_n$
\State $\zeta_n=\frac{(\bm y_n,\bm y_n)(A\bm t_n,\bm t_n)-(\bm y_n,\bm t_n)(A\bm t_n,\bm y_n)}{(A\bm t_n,A\bm t_n)(\bm y_n,\bm y_n)-(\bm y_n,A\bm t_n)(A\bm t_n,\bm y_n)}$
\State $\eta_n=\frac{(A\bm t_n,A\bm t_n)(\bm y_n,\bm t_n)-(\bm y_n,A\bm t_n)(A\bm t_n,\bm t_n)}{(A\bm t_n,A\bm t_n)(\bm y_n,\bm y_n)-(\bm y_n,A\bm t_n)(A\bm t_n,\bm y_n)}$
\If{$n=0$}
\State $\zeta_n=\frac{(A\bm t_n,\bm t_n)}{(A\bm t_n,A\bm t_n)},\quad\eta_n=0$
\EndIf
\State $\bm u_n=\zeta_nA\bm p_n+\eta_n(\bm t_{n-1}-\bm r_n+\beta_{n-1}\bm u_{n-1})$
\State $\bm z_n=\zeta_n\bm r_n+\eta_n\bm z_{n-1}-\alpha_n\bm u_n$
\State $\bm x_{n+1}=\bm x_n+\alpha_n\bm p_n+\bm z_n$
\State $\bm r_{n+1}=\bm t_n-\eta_n\bm y_n-\zeta_nA\bm t_n$
\State $\beta_n=\frac{\alpha_n}{\zeta_n}\times\frac{(\bm r_0^{*},\bm r_{n+1})}{(\bm r_0^{*},\bm r_n)}$
\State $\bm w_n=A\bm t_n+\beta_nA\bm p_n$
\EndFor\end{algorithmic}\end{algorithm}
{From Zhang's framework, the GPBiCG($m,\ell$) method was proposed by Fujino \BookCite{71}. A unified generalization of the GPBiCG method and the BiCGSTAB($\ell$) method was proposed by K. Aihara \BookCite{4}, which is referred to as the GPBiCGstab($L$) method.}
% block: product-numerical-examples page:142 type:subsection
\paragraph{{Numerical Examples}}\label{subsec:product-numerical-examples}\mbox{}\par
{In this subsection, some illustrative numerical examples are shown. We consider nonsymmetric linear systems \eqref{eq:three-dimensional-difference-system} with parameters $N=20$, $a_1=a_2=a_3=1$, $b_1=b_2=b_3=1$, and two different parameters $c=300$ and $c=900$. The stopping criterion used for all the Krylov subspace methods is $\|\bm r_n\|/\|\bm b\|\le10^{-10}$.}
{The convergence histories of the CGS method, the BiCGSTAB method, and the GPBiCG method are shown in Fig.\ \ref{fig:product-methods-moderate} for $c=300$. The horizontal axis is not the number of iterations but the number of matrix--vector multiplications. (The CGS method, the BiCGSTAB method, and the GPBiCG method require two matrix--vector multiplications per iteration.)}
{From Fig.\ \ref{fig:product-methods-moderate}, the GPBiCG method converges faster than BiCGSTAB method, and the CGS method is the best in terms of the number of matrix--vector multiplications, or equivalently the number of iterations. On the other hand, the CGS method shows irregular convergence behavior, which may lead to a loss of accuracy. The GPBiCG method and the BiCGSTAB method show relatively smooth convergence behavior, which may be caused by local minimization of the residual 2-norms. In terms of accuracy, the $\log_{10}$ of the true relative residual 2-norms for the CGS method, the BiCGSTAB method, and the GPBiCG method are $-9.79$, $-10.50$, $-10.17$ respectively. From which, all the methods produced sufficiently accurate approximate solutions.}
% block: product-residual-gap pages:142,143 type:ordinary-paragraph
{Next, the convergence histories for $c=900$ are shown in Fig.\ \ref{fig:product-methods-difficult}. We see from Fig.\ \ref{fig:product-methods-difficult} that the CGS method shows much irregular convergence behavior and the slowest convergence, and the GPBiCG method is the best in terms of the number of matrix--vector multiplications. In terms of accuracy, the $\log_{10}$ of the true relative residual 2-norms for the CGS method, the BiCGSTAB method, and the GPBiCG method are $-6.37$, $-10.03$, $-9.79$ respectively. This means that the CGS method suffers from a loss of accuracy of about four digits. This loss of accuracy is known as a residual gap. Therefore, it is strongly recommended to check the true relative residual 2-norm, i.e., $\|\bm b-A\bm x_n\|/\|\bm b\|$ (not $\|\bm r_n\|/\|\bm b\|$!), after satisfying the stopping criterion.}
% block: product-moderate-figure page:142 type:figure
\begin{figure}[H]\centering\includegraphics[width=.87\linewidth]{figures/product-methods-moderate.png}
\caption{{Convergence histories of CGS, BiCGSTAB, and GPBiCG for $N=20,c=300$}}\label{fig:product-methods-moderate}\end{figure}
% block: product-difficult-figure page:143 type:figure
\begin{figure}[H]\centering\includegraphics[width=.87\linewidth]{figures/product-methods-difficult.png}
\caption{{Convergence histories of CGS, BiCGSTAB, and GPBiCG for $N=20,c=900$}}\label{fig:product-methods-difficult}\end{figure}
% block: product-related-methods pages:143,144 type:subsection
\paragraph{{Other Related Methods}}\label{subsec:product-related-methods}\mbox{}\par
{The QMR method in Sect.\ \ref{subsec:qmr-method} requires the information of $A^{\mathrm H}$, as well as the BiCG method. Similar to the idea of the CGS method, R. W. Freund proposed a squared variant of the QMR method referred to as the TFQMR (transpose-free QMR) method \BookCite{63}, which tends to give smooth convergence behavior. Based on the BiCGSTAB method and the idea of the QMR method, the QMRCGSTAB method was proposed in \BookCite{34}.}
{The BiCR method in Sect.\ \ref{subsec:bicr-method} may give a smoother convergence behavior than the BiCG method when the coefficient matrix $A$ is close to Hermitian, since the BiCR method with the choice $\bm r_0^{*}=\bm r_0$ becomes the CR method when $A$ is Hermitian and the residual 2-norm of the CR method monotonically decreases. Thus it is natural to replace the product-type Krylov subspace methods based on the BiCG method with those based on the BiCR method. In fact, the product-type methods based on the BiCR method can be constructed, and the resulting algorithms (CRS, BiCRSTAB, GPBiCR) are derived in \BookCite{165}. K. Abe and G. L. G. Sleijpen independently proposed product-type methods based on the BiCR method, see \BookCite{2}.}
{Y.-F. Jing et al. proposed an approach similar to the BiCR method called the BiCOR method \BookCite{114}, and the product-type variants were proposed in \BookCite{32}.}
% block: idr-heading page:144 type:subsection
\subsubsection{{Induced Dimension Reduction (IDR($s$)) Method}}\label{subsec:idr-method}
{The Induced Dimension Reduction (IDR) method was developed by Wesseling and Sonneveld around 1979 \BookCite{205}. The IDR method is mathematically equivalent to the BiCGSTAB method, and the BiCGSTAB method can be regarded as a stabilized variant of the IDR method. It is written in the acknowledgements in \BookCite{195} that Sonneveld suggested that van der Vorst reconsiders Sonneveld's IDR.}
{Though the IDR method was not given further attention, about 30 years later, Sonneveld and van Gijzen proposed an innovative algorithm called the IDR($s$) method \BookCite{174}. In this section, the theory and the derivation of the IDR($s$) method are described. The derivation in this section looks to be redundant, but it makes it easier to understand the principle of the IDR($s$) method.}
% block: idr-theorem page:144 type:theorem
\par\Needspace{14\baselineskip}
\begin{Theorem}[{The IDR theorem} \BookCite{174}, \BookCite{205}, p.550]\label{thm:idr-dimension-reduction}
{Let $A$ be any matrix in $\mathbb C^{N\times N}$, let $\bm v_0$ be any nonzero vector in $\mathbb C^N$, and let $\mathcal G_0$ be the complete Krylov space, i.e., $\mathcal G_0:=\mathcal K_N(A,\bm v_0)$. Let $\mathcal S$ denote any (proper) subspace of $\mathbb C^N$ such that $\mathcal S$ and $\mathcal G_0$ do not share a nontrivial invariant subspace,}\footnote{{This means that $\mathcal S\cap\mathcal G_0$ does not contain any eigenvector of $A$, and the trivial invariant subspace is $\{\bm0\}$.}}
{and define the sequences $\mathcal G_j$, $j=1,2,\ldots$ as}
\begin{equation}\label{eq:idr-subspace-recurrence}\mathcal G_j=(I-\omega_jA)(\mathcal G_{j-1}\cap\mathcal S),\end{equation}
{where the $\omega_j$'s are nonzero scalars. Then:}
\begin{enumerate}[label=(\arabic*)]
\item $\mathcal G_j\subset\mathcal G_{j-1}$ {for all $j>0$.}
\item $\mathcal G_j=\{\bm0\}$ {for some $j\le N$.}
\end{enumerate}\end{Theorem}
% block: idr-theorem-proof-nesting pages:144,145 type:proof
\begin{Proof}
{The proof is based on \BookCite{174}, Theorem 2.1.}
{The first statement is shown by induction. First, we show $\mathcal G_1\subset\mathcal G_0$. From $\mathcal G_0=\mathcal K_N(A,\bm v_0)$, it follows that $(I-\omega_1A)\mathcal G_0\subset\mathcal G_0$, because for a given $\bm v\in(I-\omega_1A)\mathcal G_0$ we have}
\[\begin{aligned}\bm v&=(I-\omega_1A)(c_0\bm v_0+c_1A\bm v_0+\cdots+c_{N-1}A^{N-1}\bm v_0)\\&=c_0(I-\omega_1A)\bm v_0+c_1(I-\omega_1A)A\bm v_0+\cdots+c_{N-1}(I-\omega_1A)A^{N-1}\bm v_0\\&=d_0\bm v_0+d_1A\bm v_0+\cdots+d_{N-1}A^{N-1}\bm v_0+d_NA^N\bm v_0\\&=\widetilde d_0\bm v_0+\widetilde d_1A\bm v_0+\cdots+\widetilde d_{N-1}A^{N-1}\bm v_0\\&\in\mathcal G_0.\end{aligned}\]
{Thus we have $\mathcal G_1\subset\mathcal G_0$ because}
\[\mathcal G_1=(I-\omega_1A)(\mathcal G_0\cap\mathcal S)\subset(I-\omega_1A)\mathcal G_0\subset\mathcal G_0.\]
{Next we show $\mathcal G_{j+1}\subset\mathcal G_j$ (i.e., $\bm x\in\mathcal G_{j+1}\Rightarrow\bm x\in\mathcal G_j$) for all $j>0$ using the assumption $\mathcal G_j\subset\mathcal G_{j-1}$. Let $\bm x\in\mathcal G_{j+1}$. Then $\bm x$ can be written as $\bm x=(I-\omega_{j+1}A)\bm y$ for some $\bm y\in\mathcal G_j\cap\mathcal S$. From the assumption $\mathcal G_j\subset\mathcal G_{j-1}$, it follows that $\bm y\in\mathcal G_j\cap\mathcal S\subset\mathcal G_{j-1}\cap\mathcal S$, and thus $(I-\omega_jA)\bm y\in\mathcal G_j$. Since $\bm y\in\mathcal G_j$ and $\bm y-\omega_jA\bm y\in\mathcal G_j$, we have $A\bm y\in\mathcal G_j$, and thus $(I-\omega_{j+1}A)\bm y=\bm x\in\mathcal G_j$.}
% block: idr-theorem-proof-dimension page:145 type:proof
{The second statement is shown next. The property $\mathcal G_{j+1}\subset\mathcal G_j$ means that $\dim(\mathcal G_{j+1})<\dim(\mathcal G_j)$ or $\dim(\mathcal G_{j+1})=\dim(\mathcal G_j)$.}
{If $\mathcal G_j\cap\mathcal S\ne\mathcal G_j$, then $\dim(\mathcal G_j\cap\mathcal S)<\dim(\mathcal G_j)$ and thus we have $\dim(\mathcal G_{j+1})<\dim(\mathcal G_j)$ because $\dim(\mathcal G_{j+1})=\dim((I-\omega_jA)(\mathcal G_j\cap\mathcal S))\le\dim(\mathcal G_j\cap\mathcal S)<\dim(\mathcal G_j)$.}
{On the other hand, if $\mathcal G_j\cap\mathcal S=\mathcal G_j$, then $\mathcal G_j\subset\mathcal S$, and we have $\mathcal G_{j+1}=(I-\omega_{j+1}A)(\mathcal G_j\cap\mathcal S)=(I-\omega_{j+1}A)\mathcal G_j\subset\mathcal G_j$. (The last inclusion follows from the first statement $\mathcal G_{j+1}\subset\mathcal G_j$.) This implies $A\mathcal G_j\subset\mathcal G_j$, which means that $\mathcal G_j$ is an invariant subspace of $A$. Recalling $\mathcal G_j\subset\mathcal S$ and $\mathcal G_j\subset\mathcal G_0$ leads to $\mathcal G_j\subset\mathcal G_0\cap\mathcal S$. From the assumption that $\mathcal G_0\cap\mathcal S$ do not share a nontrivial invariant subspace of $A$, the subspace $\mathcal G_j$ is the trivial invariant subspace of $A$, i.e., $\mathcal G_j=\{\bm0\}$.}
{Therefore either $\dim(\mathcal G_{j+1})<\dim(\mathcal G_j)$ or $\dim(\mathcal G_j)=0$ occurs, which indicates that $\dim(\mathcal G_k)=0$, i.e., $\mathcal G_k=\{\bm0\}$ for some $k\le N$.}
\par\nopagebreak[4]
\end{Proof}
% block: idr-dimension-corollary page:145 type:corollary
\par\Needspace{3\baselineskip}
\begin{Corollary}\label{cor:idr-dimension-step}
{Let $\mathcal S$ be a subspace of $\mathbb C^N$ with $\dim(\mathcal S)=N-s$ for a given $s\ge1$. If $\mathcal G_{i-1}+\mathcal S=\mathbb C^N$ and $(I-\omega_jA)$ is nonsingular, then}
\begin{equation}\label{eq:idr-dimension-step}\dim(\mathcal G_j)=\dim(\mathcal G_{j-1})-s.\end{equation}
\end{Corollary}
% block: idr-dimension-corollary-proof page:145 type:proof
\begin{Proof}
{Let $W_1,W_2$ be subspaces of $\mathbb C^N$ (or a general vector space $V$). Then it is well known in linear algebra that $\dim(W_1)+\dim(W_2)=\dim(W_1+W_2)+\dim(W_1\cap W_2)$. Let $W_1=\mathcal G_{i-1}$, $W_2=\mathcal S$. Then, from the assumption $\mathcal G_{i-1}+\mathcal S=\mathbb C^N$, we have $\dim(\mathcal G_{i-1}+\mathcal S)=N$. Thus $\dim(\mathcal G_i)=\dim((I-\omega_jA)(\mathcal G_{i-1}\cap\mathcal S))=\dim(\mathcal G_{i-1}\cap\mathcal S)=\dim(\mathcal G_{i-1})+\dim(\mathcal S)-N=\dim(\mathcal G_{i-1})-s$.}
\par\nopagebreak[4]
\end{Proof}
% block: idr-dimension-remark page:145 type:remark
\par\Needspace{3\baselineskip}
\begin{Remark}\label{rem:idr-dimension-count}
{From Corollary\ \ref{cor:idr-dimension-step}, if $\mathcal G_{i-1}+\mathcal S=\mathbb C^N$ for $i=1,2,\ldots,n$, then $\dim(\mathcal G_n)=\dim(\mathcal G_0)-ns$. Thus, if $N/s$ is an integer and $\dim(\mathcal G_0)=N$, we have $\dim(\mathcal G_{N/s})=0$, i.e., $\mathcal G_{N/s}=\{\bm0\}$. This remark will be cited in Remark\ \ref{rem:idr-matvec-bound} for the maximum required number of matrix--vector multiplications of the IDR($s$) method.}
\end{Remark}
% block: idr-derivation-heading pages:145,146 type:subsection
\paragraph{{The Derivation of the IDR($s$) Method}}\label{subsec:idr-derivation}\mbox{}\par
{The IDR($s$) method with a given number $s\ge1$ generates residual vectors $\bm r_0,\bm r_1,\ldots$ such that:}
\begin{enumerate}
\item $\bm r_{j(s+1)},\bm r_{j(s+1)+1},\ldots,\bm r_{j(s+1)+s}\in\mathcal G_j,\quad j=0,1,2,\ldots$;
\item{The subspace $\mathcal S$ in \eqref{eq:idr-subspace-recurrence} is chosen so that $\dim\mathcal S=N-s$.}
\end{enumerate}
{For simplicity, we consider the case $s=2$, i.e., the IDR(2) method. In this case, the IDR(2) method produces the following residual vectors:}
\[\bm r_0,\bm r_1,\bm r_2\in\mathcal G_0,\quad\bm r_3,\bm r_4,\bm r_5\in\mathcal G_1,\quad\bm r_6,\bm r_7,\bm r_8\in\mathcal G_2,\ldots\]
% block: idr-two-shadow-space page:146 type:ordinary-paragraph
{Here we explain how to produce $\bm r_3,\bm r_4,\bm r_5\in\mathcal G_1$ from the given residual vectors $\bm r_0,\bm r_1,\bm r_2\in\mathcal G_0$. Since $s=2$, it follows from Remark\ \ref{rem:idr-dimension-count} that we need to choose $\mathcal S$ such that $\dim(\mathcal S)=N-2$. The standard choice is}
\begin{equation}\label{eq:idr-shadow-complement}\mathcal S=\operatorname{span}\{\bm s_1,\bm s_2\}^{\perp}\end{equation}
{with arbitrary linearly independent vectors $\bm s_1$ and $\bm s_2$, where the symbol $\perp$ is the orthogonal complement.}
{Now, let $S=[\bm s_1,\bm s_2]\in\mathbb C^{N\times2}$ with linearly independent vectors $\bm s_1$ and $\bm s_2$. Then for a vector $\bm v\in\mathbb C^N$, the following fact holds true:}
\begin{equation}\label{eq:idr-shadow-orthogonality}\bm v\perp\mathcal S\Leftrightarrow S^{\mathrm H}\bm v=\bm0.\end{equation}
{In what follows, we use the following definitions:}
\begin{equation}\label{eq:idr-residual-differences}\Delta\bm r_k:=\bm r_{k+1}-\bm r_k,\quad dR_k:=[\Delta\bm r_k,\Delta\bm r_{k+1}],\quad\bm c_k=(S^{\mathrm H}dR_k)^{-1}S^{\mathrm H}\bm r_{k+2}.\end{equation}
% block: idr-first-subspace-residuals pages:146,147 type:ordinary-paragraph
{The following vector $\bm v_0$ lies in $\mathcal G_0\cap\mathcal S$:}
\begin{equation}\label{eq:idr-first-projected-vector}\bm v_0=\bm r_2-dR_0\bm c_0\in\mathcal G_0\cap\mathcal S.\end{equation}
{In fact, $\bm v_0\in\mathcal G_0$ because $\bm v_0\in\operatorname{span}\{\bm r_2,\Delta\bm r_0,\Delta\bm r_1\}$ with $\bm r_2,\Delta\bm r_0,\Delta\bm r_1\in\mathcal G_0$ from the assumption $\bm r_0,\bm r_1,\bm r_2\in\mathcal G_0$, and $\bm v_0\in\mathcal S$ because it is easy to see $S^{\mathrm H}\bm v_0=\bm0$ and the relation \eqref{eq:idr-shadow-orthogonality}.}
{From \eqref{eq:idr-subspace-recurrence}, we can produce $\bm r_3\in\mathcal G_1$ by}
\begin{equation}\label{eq:idr-first-reduced-residual}\bm r_3=(I-\omega_0A)\bm v_0\in\mathcal G_1,\end{equation}
{where $\omega_0$ is usually chosen so that $\|\bm r_3\|$ is minimized, i.e., $\omega_0=(A\bm v_0,\bm v_0)/(A\bm v_0,A\bm v_0)$.}
{In order to produce $\bm r_4\in\mathcal G_1$, we need a vector $\bm v_1$ such that $\bm v_1\in\mathcal G_0\cap\mathcal S$. To this end, $\bm v_1$ is constructed by using $\bm r_3$ and $dR_1$ as follows:}
\[\bm v_1=\bm r_3-dR_1\bm c_1\in\mathcal G_0\cap\mathcal S.\]
{From a similar discussion to that above, we see $\bm v_1\in\mathcal G_0$ and $\bm v_1\in\mathcal S$. Thus, we have}
\[\bm r_4=(I-\omega_0A)\bm v_1\in\mathcal G_1.\]
{Similarly,}
\[\bm v_2=\bm r_4-dR_2\bm c_2\in\mathcal G_0\cap\mathcal S,\quad\bm r_5=(I-\omega_0A)\bm v_2\in\mathcal G_1.\]
% block: idr-second-subspace-residuals page:147 type:ordinary-paragraph
{Now, we obtained $\bm r_3,\bm r_4,\bm r_5\in\mathcal G_1$ from the given residual vectors $\bm r_0,\bm r_1,\bm r_2\in\mathcal G_0$. Similarly, the following algorithm produces $\bm r_6,\bm r_7,\bm r_8\in\mathcal G_2$ from $\bm r_3,\bm r_4,\bm r_5\in\mathcal G_1$:}
\[\begin{aligned}\bm v_3&=\bm r_5-dR_3\bm c_3\in\mathcal G_1\cap\mathcal S,&\bm r_6&=(I-\omega_1A)\bm v_3\in\mathcal G_2,\\\bm v_4&=\bm r_6-dR_4\bm c_4\in\mathcal G_1\cap\mathcal S,&\bm r_7&=(I-\omega_1A)\bm v_4\in\mathcal G_2,\\\bm v_5&=\bm r_7-dR_5\bm c_5\in\mathcal G_1\cap\mathcal S,&\bm r_8&=(I-\omega_1A)\bm v_5\in\mathcal G_2,\end{aligned}\]
{where $\omega_1$ is usually chosen so that $\|\bm r_6\|$ is minimized, i.e., $\omega_1=(A\bm v_3,\bm v_3)/(A\bm v_3,A\bm v_3)$.}
% block: idr-solution-extraction page:147 type:ordinary-paragraph
{Now, we describe how to extract approximate solutions $\bm x_3,\bm x_4,\bm x_5$ from $\bm x_0,\bm x_1,\bm x_2$. In what follows, we use the following definition:}
\[\Delta\bm x_k:=\bm x_{k+1}-\bm x_k,\qquad dX_k:=[\Delta\bm x_k,\Delta\bm x_{k+1}].\]
{Then from \eqref{eq:idr-first-projected-vector} and \eqref{eq:idr-first-reduced-residual} $\bm x_3$ is obtained as follows:}
\[\begin{aligned}\bm r_3&=(I-\omega_0A)\bm v_0=\bm r_2-dR_0\bm c_0-\omega_0A\bm v_0\\&\Leftrightarrow\bm b-A\bm x_3=\bm b-A\bm x_2-dR_0\bm c_0-\omega_0A\bm v_0\\&\Leftrightarrow\bm x_3=\bm x_2+A^{-1}dR_0\bm c_0+\omega_0\bm v_0\\&\Leftrightarrow\bm x_3=\bm x_2-dX_0\bm c_0+\omega_0\bm v_0.\end{aligned}\]
{Similarly,}
\[\begin{aligned}\bm x_4&=\bm x_3-dX_1\bm c_1+\omega_0\bm v_1,\\\bm x_5&=\bm x_4-dX_2\bm c_2+\omega_0\bm v_2.\end{aligned}\]
{$\bm x_6,\bm x_7,\bm x_8$ can be obtained as follows:}
\[\begin{aligned}\bm x_6&=\bm x_5-dX_3\bm c_3+\omega_1\bm v_3,\\\bm x_7&=\bm x_7-dX_4\bm c_4+\omega_1\bm v_4,\\\bm x_8&=\bm x_8-dX_5\bm c_5+\omega_1\bm v_5.\end{aligned}\]
{Finally, we describe how to determine initial residual vectors $\bm r_0,\bm r_1,\bm r_2\in\mathcal G_0$ and $\bm x_0,\bm x_1,\bm x_2$. These vectors are given as follows: we compute}
\[\begin{aligned}\bm r_0&=\bm b-A\bm x_0,\\\bm r_1&=(I-c_0A)\bm r_0,&\bm x_1&=\bm x_0+c_0\bm r_0,\\\bm r_2&=(I-c_1A)\bm r_1,&\bm x_2&=\bm x_1+c_1\bm r_1,\end{aligned}\]
{where $\bm x_0$ is an arbitrary vector, and $c_0$ and $c_1$ are chosen so that $\|\bm r_1\|$ and $\|\bm r_2\|$ are minimized.}
% block: idr-general-algorithm-introduction page:148 type:ordinary-paragraph
{From the above discussion, the IDR($s$) method with $s=2$ can be easily generalized to the IDR($s$) method with a general number $s\ge1$, and the algorithm is described in Algorithm\ \ref{alg:idr}. The symbol ``$\Delta X(:,i)$'' means that the $i$th column of matrix $\Delta X$.}
% block: idr-algorithm page:148 type:ordinary-paragraph
\begin{algorithm}[H]\caption{{The IDR($s$) method}}\label{alg:idr}\begin{algorithmic}[1]
\Statex\textbf{Input:} $\bm x_0\in\mathbb C^N,\ P\in\mathbb C^{N\times s},\ s\ge1$
\Statex\textbf{Output:} $\bm x_i$
\State $\bm r_0=\bm b-A\bm x_0$
\Statex ({The other initial residuals $\bm r_1,\ldots,\bm r_s\in\mathcal G_0$})
\For{$i=1$ {to} $s$}
\State $\bm v=A\bm r_{i-1},\quad\omega=\frac{\bm v^{\mathrm H}\bm r_{i-1}}{\bm v^{\mathrm H}\bm v}$
\State $\Delta X(:,i)=\omega\bm r_{i-1},\quad\Delta R(:,i)=-\omega\bm v$
\State $\bm x_i=\bm x_{i-1}+\Delta X(:,i);\quad\bm r_i=\bm r_{i-1}+\Delta R(:,i)$
\EndFor
\State $j=1,\quad i=s$
\State $M=P^{\mathrm H}\Delta R,\quad\bm h=P^{\mathrm H}\bm r_i$
\While{$\frac{\|\bm r_i\|}{\|\bm b\|}>\epsilon$}
\For{$k=0$ {to} $s$}
\State{Solve $\bm c$ from $M\bm c=\bm h$}
\State $\bm q=-\Delta R\bm c$
\State $\bm v=\bm r_i+\bm q$
\If{$k=0$}
\State $\bm t=A\bm v,\quad\omega=\frac{\bm t^{\mathrm H}\bm v}{\bm t^{\mathrm H}\bm t}$
\State $\Delta R(:,j)=\bm q-\omega\bm t$
\State $\Delta X(:,j)=-\Delta X\bm c+\omega\bm v$
\Else
\State $\Delta X(:,j)=-\Delta X\bm c+\omega\bm v$
\State $\Delta R(:,j)=-A\Delta X(:,j)$
\EndIf
\Statex ({Update approximate solutions $\bm x_i$})
\State $\bm r_{i+1}=\bm r_i+\Delta R(:,j)$
\State $\bm x_{i+1}=\bm x_i+\Delta X(:,j)$
\State $\Delta\bm m=P^{\mathrm H}\Delta R(:,j)$
\State $M(:,j)=\Delta\bm m$
\State $\bm h=\bm h+\Delta\bm m$
\State $i=i+1,\quad j=j+1$
\State $j=(j-1)\%s+1$ ({$\%$: modulo operation, i.e.\ $a\%n=r$, where $a=mn+r$.})
\EndFor\EndWhile\end{algorithmic}\end{algorithm}
% block: idr-matvec-remark page:148 type:remark
\par\Needspace{3\baselineskip}
\begin{Remark}\label{rem:idr-matvec-bound}
{In order to produce residual vectors in $\mathcal G_j$ from residual vectors in $\mathcal G_{j-1}$, the IDR($s$) method requires $s+1$ matrix--vector multiplications. From Remark\ \ref{rem:idr-dimension-count}, the residual vector in $\mathcal G_{N/s}$ is zero. Thus, in exact precision arithmetic, the number of matrix--vector multiplications for the IDR($s$) method is at most $(s+1)\times N/s=N+N/s$ to obtain the exact solution. Note that the number of matrix--vector multiplications for the product-type Krylov subspace methods in Sect.\ \ref{subsec:product-type-krylov} is at most $2N$.}
\end{Remark}
% block: idr-bicgstab-remark page:149 type:remark
{From the following remark, the IDR($s$) method can be regarded as an extension of the BiCGSTAB method.}
\par\Needspace{3\baselineskip}
\begin{Remark}\label{rem:idr-bicgstab-equivalence}
{The IDR(1) method and the BiCGSTAB method give the same residual vectors at the even steps.}
\end{Remark}
% block: idr-sonneveld-remark page:149 type:remark
\par\Needspace{3\baselineskip}
\begin{Remark}\label{rem:idr-sonneveld-subspace}
{In terms of orthogonality, the IDR($s$) method produces residual vectors that belong to the following so-called Sonneveld subspace:}
\begin{equation}\label{eq:idr-sonneveld-subspace}\{\Omega_j(A)\bm v:\bm v\perp\mathcal B_j(A^{\mathrm H},S)\}.\end{equation}
{Here $\Omega_j(A):=\prod_{i=1}^j(I-\omega_iA)$, and $\mathcal B_j(A^{\mathrm H},S)$ is a subspace of $\mathbb C^N$ generated by using $A$ and $S:=[\bm s_1,\ldots,\bm s_s]$, which is defined by}
\[\mathcal B_j(A^{\mathrm H},S):=\mathcal K_j(A^{\mathrm H},\bm s_1)+\cdots+\mathcal K_j(A^{\mathrm H},\bm s_s).\]
{Here the symbol ``$+$'' means the sum of subspaces.}\footnote{{Given two subspaces $W_1$ and $W_2$, the sum of $W_1$ and $W_2$ is defined by $W_1+W_2:=\{\bm w_1+\bm w_2:\bm w_1\in W_1,\bm w_2\in W_2\}$.}}
{$\mathcal B_j(A^{\mathrm H},R_0^{*})$ will be redefined in Definition\ \ref{def:block-krylov-subspace} for block Krylov subspace methods. For the details of the view in \eqref{eq:idr-sonneveld-subspace}, see \BookCite{161}, \BookCite{163}.}
\end{Remark}
% block: idr-sonneveld-first-level page:149 type:ordinary-paragraph
{We give a rough explanation of Remark\ \ref{rem:idr-sonneveld-subspace} using the IDR($s$) method with $s=2$ for the two cases $j=1,2$ in \eqref{eq:idr-sonneveld-subspace}. First we consider the case $j=1$. From \eqref{eq:idr-subspace-recurrence} $\mathcal G_1$ is generated by}
\[\mathcal G_1=(I-\omega_1A)(\mathcal G_0\cap\mathcal S),\]
{where $\mathcal S=\operatorname{span}\{\bm s_1,\bm s_2\}^{\perp}$, see also \eqref{eq:idr-shadow-complement}. Then, $\bm r\in\mathcal G_1$ can be written as}
\[\bm r=(I-\omega_1A)\bm v,\quad\bm v\in\mathcal G_0,\quad\bm v\perp\operatorname{span}\{\bm s_1,\bm s_2\}.\]
{Note that $\bm v\in\mathcal S$ means $\bm v\perp\operatorname{span}\{\bm s_1,\bm s_2\}$. Since $\mathcal B_1(A^{\mathrm H},S)=\mathcal K_1(A^{\mathrm H},\bm s_1)+\mathcal K_1(A^{\mathrm H},\bm s_2)=\operatorname{span}\{\bm s_1,\bm s_2\}$, we have}
\[\bm r\in\{(I-\omega_1A)\bm v:\bm v\perp\mathcal B_1(A^{\mathrm H},S)\}=\{\Omega_1(A)\bm v:\bm v\perp\mathcal B_1(A^{\mathrm H},S)\},\]
{which corresponds to the case $j=1$ in \eqref{eq:idr-sonneveld-subspace}.}
% block: idr-sonneveld-second-level pages:149,150 type:ordinary-paragraph
{Next, we consider the case $j=2$. $\mathcal G_2$ is generated by}
\[\mathcal G_2=(I-\omega_2A)(\mathcal G_1\cap\mathcal S).\]
{Then a residual vector $\widetilde{\bm r}\in\mathcal G_2$ can be written as}
\begin{equation}\label{eq:idr-second-level-residual}\widetilde{\bm r}=(I-\omega_2A)\widetilde{\bm v},\quad\widetilde{\bm v}\in\mathcal G_1,\quad\widetilde{\bm v}\perp\operatorname{span}\{\bm s_1,\bm s_2\}.\end{equation}
{Since $\widetilde{\bm v}\in\mathcal G_1=(I-\omega_1A)(\mathcal G_0\cap\mathcal S)$, the vector $\widetilde{\bm v}$ can be written as}
\begin{equation}\label{eq:idr-first-level-preimage}\widetilde{\bm v}=(I-\omega_1A)\bm v,\quad\bm v\in\mathcal G_0,\quad\bm v\perp\operatorname{span}\{\bm s_1,\bm s_2\}.\end{equation}
{From \eqref{eq:idr-second-level-residual} and \eqref{eq:idr-first-level-preimage}, we have}
\begin{equation}\label{eq:idr-two-factor-residual}\widetilde{\bm r}=(I-\omega_2A)(I-\omega_1A)\bm v,\end{equation}
{where}
\[\widetilde{\bm v}=(I-\omega_1A)\bm v\perp\operatorname{span}\{\bm s_1,\bm s_2\},\quad\bm v\perp\operatorname{span}\{\bm s_1,\bm s_2\}.\]
{Note that}
\[\begin{aligned}(I-\omega_1A)\bm v\perp\operatorname{span}\{\bm s_1,\bm s_2\}&\Leftrightarrow((I-\omega_1A)\bm v)^{\mathrm H}\bm s_i=0\quad(i=1,2)\\&\Leftrightarrow\bm v^{\mathrm H}(I-\omega_1A)^{\mathrm H}\bm s_i=0\quad(i=1,2)\\&\Leftrightarrow\bm v\perp\operatorname{span}\{(I-\omega_1A)^{\mathrm H}\bm s_1,(I-\omega_1A)^{\mathrm H}\bm s_2\}.\end{aligned}\]
{$\bm v\perp\operatorname{span}\{(I-\omega_1A)^{\mathrm H}\bm s_1,(I-\omega_1A)^{\mathrm H}\bm s_2\}$ and $\bm v\perp\operatorname{span}\{\bm s_1,\bm s_2\}$ imply}
\[\bm v\perp\bm s_1,\quad\bm v\perp\bm s_2,\quad\bm v\perp A^{\mathrm H}\bm s_1,\quad\bm v\perp A^{\mathrm H}\bm s_1.\]
{Thus,}
\[\begin{aligned}\bm v\perp\operatorname{span}\{\bm s_1,A^{\mathrm H}\bm s_1,\bm s_2,A^{\mathrm H}\bm s_2\}.\\=\operatorname{span}\{\bm s_1,A^{\mathrm H}\bm s_1\}+\operatorname{span}\{\bm s_2,A^{\mathrm H}\bm s_2\}\\=\mathcal K_2(A^{\mathrm H},\bm s_1)+\mathcal K_2(A^{\mathrm H},\bm s_2)\\=\mathcal B_2(A^{\mathrm H},S).\end{aligned}\]
{Then, from \eqref{eq:idr-two-factor-residual}, together with the above result, we have}
\begin{equation}\label{eq:idr-second-level-block-orthogonality}\widetilde{\bm r}=(I-\omega_2A)(I-\omega_1A)\bm v,\quad\bm v\perp\mathcal B_2(A^{\mathrm H},S),\end{equation}
{from which we obtain}
\begin{equation}\label{eq:idr-second-sonneveld-membership}\widetilde{\bm r}\in\{\Omega_2(A)\bm v:\bm v\perp\mathcal B_2(A^{\mathrm H},S)\}.\end{equation}
{This corresponds to the case $j=2$ in \eqref{eq:idr-sonneveld-subspace}.}
% block: idr-related-literature pages:150,151 type:ordinary-paragraph
{In what follows, some references related to the IDR($s$) method are described. An elegant explanation of the IDR($s$) method is found in \BookCite{88}, where a visualization of the IDR theorem in \BookCite{88}, Fig.3.1 is useful for intuitively understanding the IDR theorem (Theorem\ \ref{thm:idr-dimension-reduction}). A more stable and accurate variant of the IDR($s$) method is proposed by the same authors in \BookCite{199}, and the IDR($s$) method with the quasi-minimal residual strategy is proposed in \BookCite{47}. Variants of the IDR($s$) method with partial orthonormalization is proposed in \BookCite{209}. Extensions of the IDR($s$) methods using the idea of the BiCGSTAB($\ell$) method are known as the IDRstab method \BookCite{164} and the GBiCGSTAB($s,L$) method \BookCite{185}. Variants of the IDRstab method are found in \BookCite{5}--\BookCite{7}.}
% block: idr-numerical-example page:151 type:subsection
\paragraph{{Numerical Example}}\label{subsec:idr-numerical-example}\mbox{}\par
{In this subsection, typical convergence behavior of the IDR($s$) method is shown by one numerical experiment. We consider nonsymmetric linear systems \eqref{eq:three-dimensional-difference-system} with parameters $N=20$, $a_1=a_2=a_3=1$, $b_1=b_2=b_3=1$, and $c=300$. The stopping criterion is $\|\bm r_n\|/\|\bm b\|\le10^{-10}$.}
{The convergence histories of the IDR($s$) method ($s=1,2,4$) and the BiCGSTAB method are shown in Fig.\ \ref{fig:idr-comparison}. The horizontal axis is the number of matrix--vector multiplications.}
{From Fig.\ \ref{fig:idr-comparison}, we see that as the number of $s$ in the IDR($s$) method increases, the required number of matrix--vector multiplications decreases. The IDR(4) method and the IDR(2) method converge faster than the BiCGSTAB method. In terms of accuracy, the true relative residual 2-norms for the BiCGSTAB method, the IDR(1) method, the IDR(2) method, and the IDR(4) method are $-10.50$, $-10.04$, $-10.2$, $-8.76$.}
% block: idr-comparison-figure page:151 type:figure
\begin{figure}[H]\centering\includegraphics[width=.87\linewidth]{figures/idr-comparison.png}
\caption{{Convergence histories of BiCGSTAB, IDR(1), IDR(2), and IDR(4)}}\label{fig:idr-comparison}\end{figure}
% block: block-idr-heading page:152 type:subsection
\subsubsection{{Block Induced Dimension Reduction (Block IDR($s$)) Method}}\label{subsec:block-idr-method}
{The block IDR($s$) method \BookCite{46} is an extension of the IDR($s$) method to solve the block linear systems of the form}
\begin{equation}\label{eq:block-linear-system}AX=B,\end{equation}
{where $A\in\mathbb C^{N\times N}$, $B\in\mathbb C^{N\times m}$, and $m$ is usually much less than $N$.}
{Prior to the theory of the block IDR($s$) method, the framework of block Krylov subspace methods is briefly described after the following definitions:}
% block: block-krylov-definition page:152 type:definition
\par\Needspace{3\baselineskip}
\begin{Definition}{{Block Krylov subspace}}\label{def:block-krylov-subspace}
{Let $V\in\mathbb C^{N\times m}$. Then $\mathcal K_n(A,V)\subset\mathbb C^{N\times m}$ defined by}
\[\begin{aligned}\mathcal K_n(A,V)&:=\{VC_0+AVC_1+\cdots+A^{n-1}VC_{n-1}:C_0,C_1,\ldots,C_{n-1}\in\mathbb C^{m\times m}\}\\&(=\operatorname{block\ span}\{V,AV,\ldots,A^{n-1}V\})\end{aligned}\]
{is called the block Krylov subspace with respect to $A$ and $V$.}
\end{Definition}
% block: block-krylov-sum-definition page:152 type:definition
\par\Needspace{3\baselineskip}
\begin{Definition}{{Sum of Krylov subspaces}}\label{def:block-krylov-sum}
{Let $V=[\bm v_1,\bm v_2,\ldots,\bm v_m]\in\mathbb C^{N\times m}$. Then $\mathcal B_n(A,V)\subset\mathbb C^N$ is defined by}
\[\begin{aligned}\mathcal B_n(A,V)&:=\mathcal K_n(A,\bm v_1)+\mathcal K_n(A,\bm v_2)+\cdots+\mathcal K_n(A,\bm v_m)\\&=\{V\bm c_0+AV\bm c_1+\cdots+A^{n-1}V\bm c_{n-1}:\bm c_0,\bm c_1,\ldots,\bm c_{n-1}\in\mathbb C^m\}.\end{aligned}\]
\end{Definition}
% block: block-krylov-grade-definition page:152 type:definition
\par\Needspace{3\baselineskip}
\begin{Definition}{{Grade of block Krylov subspace}}\label{def:block-krylov-grade}
{The smallest number $n$ such that}
\[\dim(\mathcal B_n(A,V))=\dim(\mathcal B_{n+1}(A,V))\]
{is called the block grade of $A$ with respect to $V$, which is denoted by $\nu(V,A)$.}
\end{Definition}
% block: block-grade-reduction-remark page:152 type:remark
\par\Needspace{3\baselineskip}
\begin{Remark}\label{rem:block-grade-scalar-reduction}
{Definitions\ \ref{def:block-krylov-subspace} and \ref{def:block-krylov-grade} are extensions of Definitions\ \ref{def:krylov-subspace} and \ref{def:krylov-grade}. In fact, if $m=1$, these definitions are identical to Definitions\ \ref{def:krylov-subspace} and \ref{def:krylov-grade}.}
\end{Remark}
% block: block-krylov-framework pages:152,153 type:ordinary-paragraph
{We now describe the framework of the block Krylov subspace methods. Let $X_0\in\mathbb C^{N\times s}$ be an initial guess, and let $R_0:=B-AX_0$ be the corresponding initial residual matrix. Then, block Krylov subspace methods find the $n$th approximate solution $X_n$ as follows:}
\[X_n=X_0+Z_n,\quad Z_n\in\mathcal K_n(A,R_0).\]
{The corresponding residual matrix is given by}
\[R_n=B-AX_n=R_0-AZ_n\in\mathcal K_{n+1}(A,R_0).\]
{Determining $Z_n$ yields various block Krylov subspace methods, e.g.,}
\begin{itemize}
\item{The block CG method:} $R_n\perp\mathcal B_n(A,R_0)$.
\item{The block COCG method:} $R_n\perp\overline{\mathcal B_n(A,R_0)}$.
\item{The block BiCG method:} $R_n\perp\mathcal B_n(A^{\mathrm H},R_0^{*})$.
\end{itemize}
{Here, $R_0^{*}$ is an initial shadow residual matrix,}\footnote{{In practice, $R_0^{*}$ is usually set to $R_0^{*}=R_0$ or a random matrix.}}
{and the symbol $R_n\perp\mathcal B_n(A,V)$ means that all the column vectors of $R_n$ are orthogonal to $\mathcal B_n(A,V)$.}
% block: block-krylov-history page:153 type:ordinary-paragraph
{Here, we give a brief history of the block Krylov subspace methods. In 1980, O’Leary proposed the block BiCG method (and the block CG method) in \BookCite{142}. In 1990, Vital proposed the block GMRES method in \BookCite{202}. Block Krylov subspace methods based on the QMR method were proposed in 1997 \BookCite{65}, \BookCite{157}. In 2003, Guennouni et al. proposed the block BiCGSTAB method \BookCite{86}. In 2011, Du et al. proposed the block IDR($s$) method \BookCite{46}. For other block Krylov subspace methods, see \BookCite{153}, \BookCite{181}, \BookCite{182}, \BookCite{210}.}
{Now we turn our attention to deriving the block IDR($s$) method. The block IDR($s$) method is based on the following theorem (see \BookCite{46}), which is a slight modification of Theorem\ \ref{thm:idr-dimension-reduction}.}
% block: block-idr-theorem page:153 type:theorem
\par\Needspace{14\baselineskip}
\begin{Theorem}\label{thm:block-idr-dimension-reduction}
{Let $A$ be any matrix in $\mathbb C^{N\times N}$, let $\bm v_0$ be any nonzero vector in $\mathbb C^N$, and let $\mathcal G_0=\mathcal B_{\nu(R_0,A)}(A,R_0)$. Let $\mathcal S$ denote any (proper) subspace of $\mathbb C^N$ such that $\mathcal S$ and $\mathcal G_0$ do not share a nontrivial invariant subspace,}\footnote{{This means that $\mathcal S\cap\mathcal G_0$ does not contain any eigenvector of $A$, and the trivial invariant subspace is $\{\bm0\}$.}}
{and define the sequences $\mathcal G_j$, $j=1,2,\ldots$ as}
\begin{equation}\label{eq:block-idr-subspace-recurrence}\mathcal G_j=(I-\omega_jA)(\mathcal G_{j-1}\cap\mathcal S),\end{equation}
{where the $\omega_j$'s are nonzero scalars. Then:}
\begin{enumerate}[label=(\arabic*)]
\item $\mathcal G_j\subset\mathcal G_{j-1}$ {for all $j>0$.}
\item $\mathcal G_j=\{\bm0\}$ {for some $j\le N$.}
\end{enumerate}\end{Theorem}
% block: block-idr-theorem-proof pages:153,154 type:proof
\begin{Proof}
{The only difference between Theorems\ \ref{thm:idr-dimension-reduction} and \ref{thm:block-idr-dimension-reduction} is the definition of $\mathcal G_0$, i.e., $\mathcal G_0=\mathcal K_N(A,\bm v_0)$ or $\mathcal G_0=\mathcal B_{\nu(R_0,A)}(A,R_0)$. In the proof of Theorem\ \ref{thm:idr-dimension-reduction}-(1), $\mathcal G_0=\mathcal K_N(A,\bm v_0)$ is used only for proving $\mathcal G_1\subset\mathcal G_0$. Thus, all we have to do is to prove $\mathcal G_1\subset\mathcal G_0$. From $\mathcal G_0=\mathcal B_{\nu(R_0,A)}(A,R_0)$, we have $(I-\omega_1A)\mathcal G_0\subset\mathcal G_0$, because for a given $\bm v\in(I-\omega_1A)\mathcal G_0=(I-\omega_1A)\mathcal B_{\nu(R_0,A)}(A,R_0)$ we have}
\[\begin{aligned}\bm v&=(I-\omega_1A)(R_0\bm c_0+AR_0\bm c_1+\cdots+A^{\nu(R_0,A)-1}R_0\bm c_{\nu(R_0,A)-1})\\&=(I-\omega_1A)R_0\bm c_0+(I-\omega_1A)AR_0\bm c_1+\cdots+(I-\omega_1A)A^{\nu(R_0,A)-1}R_0\bm c_{\nu(R_0,A)-1}\\&=R_0\bm d_0+AR_0\bm d_1+\cdots+A^{\nu(R_0,A)-1}R_0\bm d_{\nu(R_0,A)-1}+A^{\nu(R_0,A)}R_0\bm d_{\nu(R_0,A)}\\&=R_0\widetilde{\bm d}_0+AR_0\widetilde{\bm d}_1+\cdots+A^{\nu(R_0,A)-1}R_0\widetilde{\bm d}_{\nu(R_0,A)-1}\\&\in\mathcal G_0.\end{aligned}\]
{Thus $\mathcal G_1=(I-\omega_1A)(\mathcal G_0\cap\mathcal S)\subset(I-\omega_1A)\mathcal G_0\subset\mathcal G_0$.}
{Theorem\ \ref{thm:block-idr-dimension-reduction}-(2) is shown by following the proof of Theorem\ \ref{thm:idr-dimension-reduction}-(2) and using $\dim\mathcal G_0\le N$.}
\par\nopagebreak[4]
\end{Proof}
% block: block-idr-dimension-corollary page:154 type:corollary
{Similar to Corollary\ \ref{cor:idr-dimension-step}, we have the following fact:}
\par\Needspace{3\baselineskip}
\begin{Corollary}\label{cor:block-idr-dimension-step}
{Let $\mathcal S$ be a subspace of $\mathbb C^N$ with $\dim(\mathcal S)=N-sm$ for a given $s,m\ge1$. If $\mathcal G_{i-1}+\mathcal S=\mathbb C^N$ and $(I-\omega_jA)$ is nonsingular, then}
\[\dim(\mathcal G_j)=\dim(\mathcal G_{j-1})-sm.\]
\end{Corollary}
{The difference between Corollaries\ \ref{cor:idr-dimension-step}, \ref{cor:block-idr-dimension-step} is the parameter $m$. The parameter $m$ corresponds to the size of the number of right-hand sides of $AX=B$, where $B$ is an $N$-by-$m$ matrix.}
% block: block-idr-subspace-choice page:154 type:ordinary-paragraph
{In what follows, the block IDR($s$) method is explained. The block IDR($s$) method generates residual matrices $R_k=B-AX_k$ for $k=0,1,\ldots$ such that:}
\begin{enumerate}
\item{all the column vectors of $R_{j(s+1)},R_{j(s+1)+1},\ldots,R_{j(s+1)+s}$ belong to $\mathcal G_j$ in \eqref{eq:block-idr-subspace-recurrence} for $j=0,1,2,\ldots$;}
\item{the subspace $\mathcal S$ in \eqref{eq:block-idr-subspace-recurrence} is chosen so that $\dim\mathcal S=N-sm$.}
\end{enumerate}
{Let $S=[\bm s_1,\bm s_2,\ldots,\bm s_{sm}]\in\mathbb C^{N\times sm}$. Then, a standard choice of the subspace $\mathcal S$ is as follows:}
\[\mathcal S=\operatorname{span}\{S\}^{\perp}=\operatorname{span}\{\bm s_1,\bm s_2,\ldots,\bm s_{sm}\}^{\perp},\]
{where all the column vectors of $S\in\mathbb C^{N\times sm}$ are linearly independent, and thus it is easy to see that $\dim\mathcal S=N-sm$.}
% block: block-idr-two-residuals page:154 type:ordinary-paragraph
{We now briefly explain the block IDR($s$) method with $s=2$. For $j=0,1,2$, the block IDR(2) method produces the following residual matrices:}
\[R_0,R_1,R_2,\quad R_3,R_4,R_5,\quad R_6,R_7,R_8,\]
{where all the column vectors $R_0,R_1,R_2$ belong to $\mathcal G_0$, all the column vectors $R_3,R_4,R_5$ belong to $\mathcal G_1$, and all the column vectors $R_3,R_4,R_5$ belong to $\mathcal G_2$.}
{In what follows, we describe how to obtain $R_3,R_4,R_5$ from $R_0,R_1,R_2$ under the assumption that all the column vectors of $R_0,R_1,R_2$ belong to $\mathcal G_0$. Similar to the IDR($s$) method, the following definitions are useful for the block IDR($s$) method:}
\begin{equation}\label{eq:block-idr-residual-differences}\Delta R_k:=R_{k+1}-R_k,\quad dR_k^{\blacksquare}:=[\Delta R_k,\Delta R_{k+1}],\quad C_k=(S^{\mathrm H}dR_k^{\blacksquare})^{-1}S^{\mathrm H}R_{k+2}.\end{equation}
{All the column vectors of the following matrix $V_0$ lie in $\mathcal G_0\cap\mathcal S$:}
\[V_0=R_2-dR_0^{\blacksquare}C_0,\quad\operatorname{Im}(V_0)\subset\mathcal G_0\cap\mathcal S\]
% block: block-idr-first-reduced-residual page:155 type:ordinary-paragraph
{because all the column vectors of $R_2\in\mathbb C^{N\times m}$ and $dR_0\in\mathbb C^{N\times2m}$ belong to $\mathcal G_0$ from the assumption that all the column vectors of $R_0,R_1,R_2$ belong to $\mathcal G_0$, and all the column vectors of $V_0$ also belong to $\mathcal S$ from the relation $S^{\mathrm H}V_0=O$, where $O$ is the $2m\times m$ zero matrix. Here $\operatorname{Im}(V_0)$ is the image of $V_0$, i.e., the subspace spanned by all the column vectors of $V_0$.}
{From \eqref{eq:block-idr-subspace-recurrence}, the residual matrix $R_3$ whose column vectors belong to $\mathcal G_1$ can be produced by}
\begin{equation}\label{eq:block-idr-first-reduced-residual}R_3=(I-\omega_0A)V_0,\quad\operatorname{Im}(R_3)\subset\mathcal G_1,\end{equation}
{where $\omega_0$ is usually chosen so that $\|R_3\|_{\mathrm F}$ is minimized,}\footnote{{For $R=(r_{i,j})\in\mathbb C^{N\times m}$, the symbol $\|R\|_{\mathrm F}:=(\sum_{i=1}^N\sum_{j=1}^m\bar r_{i,j}r_{i,j})^{1/2}$ is called the Frobenius norm of $R$, and for a square matrix $A=(a_{i,j})\in\mathbb C^{N\times N}$, the symbol $\operatorname{Tr}(A):=\sum_{i=1}^Na_{i,i}$ is called the trace of $A$.}}
{i.e.,}
\[\omega_0=\frac{\operatorname{Tr}(T_0^{\mathrm H}V_0)}{\operatorname{Tr}(T_0^{\mathrm H}T_0)},\quad T_0:=AV_0.\]
% block: frobenius-minimization-proposition page:155 type:proposition
{The Frobenius norm minimization follows from the following fact:}
\par\Needspace{3\baselineskip}
\begin{Proposition}\label{prop:frobenius-scalar-minimization}
{Let $V,W\in\mathbb C^{m\times n}$. Then the following minimization problem:}
\[\min_{\omega\in\mathbb C}\|V-\omega W\|_{\mathrm F}\]
{can be solved by}
\[\omega=\frac{\operatorname{Tr}(W^{\mathrm H}V)}{\operatorname{Tr}(W^{\mathrm H}W)}.\]
\end{Proposition}
% block: frobenius-minimization-proof page:155 type:proof
\begin{Proof}
{Let $V=[\bm v_1,\ldots,\bm v_n]$ and $W=[\bm w_1,\ldots,\bm w_n]$, then using}
\[\bm v:=\operatorname{vec}(V)=\begin{bmatrix}\bm v_1\\\vdots\\\bm v_n\end{bmatrix},\quad\bm w:=\operatorname{vec}(W)=\begin{bmatrix}\bm w_1\\\vdots\\\bm w_n\end{bmatrix}\]
{yields $\|V-\omega W\|_{\mathrm F}=\|\bm v-\omega\bm w\|$. Thus from \eqref{eq:scalar-residual-minimizer}, the minimizer is}
\[\omega=\frac{\bm w^{\mathrm H}\bm v}{\bm w^{\mathrm H}\bm w}=\frac{\operatorname{Tr}(W^{\mathrm H}V)}{\operatorname{Tr}(W^{\mathrm H}W)},\]
{which concludes the proof.}
\par\nopagebreak[4]
\end{Proof}
% block: block-idr-next-residuals pages:155,156 type:ordinary-paragraph
{Similar to the derivation of $\bm r_3$ and $\bm r_4$ in the IDR($s$) method, $R_3$ and $R_4$ are obtained as follows:}
\[\begin{aligned}V_1&=R_3-dR_1^{\blacksquare}C_1,&\operatorname{Im}(V_1)&\subset\mathcal G_0\cap\mathcal S,\\R_4&=(I-\omega_0A)V_1,&\operatorname{Im}(R_4)&\subset\mathcal G_1,\\V_2&=R_4-dR_2^{\blacksquare}C_2,&\operatorname{Im}(V_2)&\subset\mathcal G_0\cap\mathcal S,\\R_5&=(I-\omega_0A)V_2,&\operatorname{Im}(R_5)&\subset\mathcal G_1.\end{aligned}\]
{From the above calculations, we have residual matrices $R_k$ ($k=0,\ldots,5$) such that}
\[\begin{aligned}\operatorname{Im}(R_k)&\subset\mathcal G_0\quad\text{for }k=0,1,2,\\\operatorname{Im}(R_k)&\subset\mathcal G_1\quad\text{for }k=3,4,5.\end{aligned}\]
% block: block-idr-solution-extraction page:156 type:ordinary-paragraph
{Next, we describe how to obtain approximate solution matrices $X_3,X_4,X_5$ from $X_0,X_1,X_2$. Let}
\[\Delta X_k:=X_{k+1}-X_k,\qquad dX_k^{\blacksquare}:=[\Delta X_k,\Delta X_{k+1}].\]
{Then $X_3$ is obtained as follows:}
\[\begin{aligned}R_3&=(I-\omega_0A)V_0=R_2-dR_0^{\blacksquare}C_0-\omega_0AV_0\\&\Leftrightarrow B-AX_3=B-AX_2-dR_0^{\blacksquare}C_0-\omega_0AV_0\\&\Leftrightarrow X_3=X_2+A^{-1}dR_0^{\blacksquare}C_0+\omega_0V_0\\&\Leftrightarrow X_3=X_2-dX_0^{\blacksquare}C_0+\omega_0V_0.\end{aligned}\]
{Similarly,}
\[\begin{aligned}X_4&=X_3-dX_1^{\blacksquare}C_1+\omega_0V_1,\\X_5&=X_4-dX_2^{\blacksquare}C_2+\omega_0V_2.\end{aligned}\]
{We have described how to obtain $R_k$ and $X_k$ for $k=3,4,5$ from $R_k$ and $X_k$ for $k=0,1,2$ in the block IDR($s$) method with $s=2$. Following the above derivation, we have the block IDR($s$) method in Algorithm\ \ref{alg:block-idr}. Further, the preconditioned version of the block IDR($s$) method is shown in Algorithm\ \ref{alg:preconditioned-block-idr}. Note that the symbol ``$B(:,\ell)$'' represents the $\ell$th column of matrix $B$, and the symbol ``$\Delta X(:,im+1:(i+1)m)$'' represents the $(im+1)$th, $(im+2)$th, $\ldots$, $(i+1)m$th columns of matrix $\Delta X$.}
% block: block-idr-cost-remark pages:156,157 type:remark
\par\Needspace{3\baselineskip}
\begin{Remark}\label{rem:block-idr-matvec-bound}
{From Corollary\ \ref{cor:block-idr-dimension-step}, the dimension reduction is usually $sm$, i.e., $\dim(\mathcal G_j)=\dim(\mathcal G_{j-1})-sm$ and the block IDR($s$) method requires $m(s+1)$ matrix--vector multiplications. Thus, in exact precision arithmetic, the total number of matrix--vector multiplications for the block IDR($s$) method is at most $m(s+1)\times N/sm=N+N/s$ to obtain the exact solution matrix, which does not depend on the block size $m$. Note that if the IDR($s$) method is applied to $AX=B$ (i.e., $A\bm x_1=\bm b_1,A\bm x_2=\bm b_2,\ldots,A\bm x_m=\bm b_m$), then the required number of matrix--vector multiplications of the IDR($s$) method is at most $m(N+N/s)$, which is $m$ times larger than that of the block IDR($s$) method.}
\end{Remark}
% block: block-idr-bicgstab-remark page:157 type:remark
{From the following remark, the block IDR($s$) method can be regarded as an extension of the block BiCGSTAB method.}
\par\Needspace{3\baselineskip}
\begin{Remark}\label{rem:block-idr-bicgstab-equivalence}
{The block IDR(1) method and the block BiCGSTAB method give the same residual vectors at the even steps.}
\end{Remark}
% block: block-idr-algorithm page:157 type:ordinary-paragraph
\begin{algorithm}[H]\caption{{The block IDR($s$) method}}\label{alg:block-idr}\begin{algorithmic}[1]
\Statex\textbf{Input:} $X_0\in\mathbb C^{N\times s},\ R_0=B-AX_0,\ P\in\mathbb C^{n\times sm}$
\Statex\textbf{Output:} $X_i$
\For{$i=0$ {to} $s-1$}
\State $V=AR_i,\quad\omega=\frac{\operatorname{Tr}(V^{\mathrm H}R_i)}{\operatorname{Tr}(V^{\mathrm H}V)}$
\State $\Delta X(:,im+1:(i+1)m)=\omega R_i,\quad\Delta R(:,im+1:(i+1)m)=-\omega V$
\State $X_{i+1}=X_i+\Delta X(:,im+1:(i+1)m),\quad R_{i+1}=R_i+\Delta R(:,im+1:(i+1)m)$
\EndFor
\State $j=1,\quad i=s$
\State $M=P^{\mathrm H}\Delta R,\quad H=P^{\mathrm H}R_i$
\While{$\max_{\ell\in\{1,2,\ldots,m\}}\frac{\|R_i(:,\ell)\|_{\mathrm F}}{\|B(:,\ell)\|_{\mathrm F}}>\epsilon$}
\For{$k=0$ {to} $s$}
\State{Solve $C$ from $MC=H$}
\State $Q=-\Delta RC,$
\State $V=R_i+Q$
\If{$k=0$}
\State $T=AV,\quad\omega=\frac{\operatorname{Tr}(T^{\mathrm H}V)}{\operatorname{Tr}(T^{\mathrm H}T)}$
\State $\Delta R(:,(j-1)m+1:jm)=Q-\omega T$
\State $\Delta X(:,(j-1)m+1:jm)=-\Delta XC+\omega V$
\Else
\State $\Delta X(:,(j-1)m+1:jm)=-\Delta XC+\omega V$
\State $\Delta R(:,(j-1)m+1:jm)=-A\Delta X(:,(j-1)m+1:jm)$
\EndIf
\State $R_{i+1}=R_i+\Delta R(:,(j-1)m+1:jm)$
\State $X_{i+1}=X_i+\Delta X(:,(j-1)m+1:jm)$
\State $\Delta M=P^{\mathrm H}\Delta R(:,(j-1)m+1:jm)$
\State $M(:,(j-1)m+1:jm)=\Delta M$
\State $H=H+\Delta M$
\State $i=i+1,\quad j=j+1$
\State $j=(j-1)\%s+1$ ({$\%$: modulo operation, i.e.\ $a\%n=r$, where $a=mn+r$.})
\EndFor\EndWhile\end{algorithmic}\end{algorithm}
% block: preconditioned-block-idr-algorithm page:158 type:ordinary-paragraph
\begin{algorithm}[H]\caption{{The preconditioned block IDR($s$) method}}\label{alg:preconditioned-block-idr}\begin{algorithmic}[1]
\Statex\textbf{Input:} $X_0\in\mathbb C^{N\times s},\ R_0=B-AX_0,\ P\in\mathbb C^{n\times sm}$
\Statex\textbf{Input:} $X_0=KX_0$ ($K$: {preconditioning matrix})
\Statex\textbf{Output:} $X_i$
\For{$i=0$ {to} $s-1$}
\State $W=K^{-1}R_i$
\State $V=AW,\quad\omega=\frac{\operatorname{Tr}((K_1^{-1}V)^{\mathrm H}K_1^{-1}R_i)}{\operatorname{Tr}((K_1^{-1}V)^{\mathrm H}K_1^{-1}V)}$
\State $\Delta X(:,im+1:(i+1)m)=\omega R_i,\quad\Delta R(:,im+1:(i+1)m)=-\omega V$
\State $X_{i+1}=X_i+\Delta X(:,im+1:(i+1)m),\quad R_{i+1}=R_i+\Delta R(:,im+1:(i+1)m)$
\EndFor
\State $j=1,\quad i=s$
\State $M=P^{\mathrm H}\Delta R,\quad H=P^{\mathrm H}R_i$
\While{$\max_{\ell\in\{1,2,\ldots,m\}}\frac{\|R_i(:,\ell)\|_{\mathrm F}}{\|B(:,\ell)\|_{\mathrm F}}>\epsilon$}
\For{$k=0$ {to} $s$}
\State{Solve $C$ from $MC=H$}
\State $Q=-\Delta RC$
\State $V=R_i+Q$
\If{$k=0$}
\State $T=AK^{-1}V,\quad\omega=\frac{\operatorname{Tr}((K_1^{-1}T)^{\mathrm H}K_1^{-1}V)}{\operatorname{Tr}((K_1^{-1}T)^{\mathrm H}K_1^{-1}T)}$
\State $\Delta R(:,(j-1)m+1:jm)=Q-\omega T$
\State $\Delta X(:,(j-1)m+1:jm)=-\Delta XC+\omega V$
\Else
\State $\Delta X(:,(j-1)m+1:jm)=-\Delta XC+\omega V$
\State $\Delta R(:,(j-1)m+1:jm)=-AK^{-1}\Delta X(:,(j-1)m+1:jm)$
\EndIf
\State $R_{i+1}=R_i+\Delta R(:,(j-1)m+1:jm)$
\State $X_{i+1}=X_i+\Delta X(:,(j-1)m+1:jm)$
\State $\Delta M=P^{\mathrm H}\Delta R(:,(j-1)m+1:jm)$
\State $M(:,(j-1)m+1:jm)=\Delta M$
\State $H=H+\Delta M$
\State $i=i+1,\quad j=j+1$
\State $j=(j-1)\%s+1$ ({$\%$: modulo operation, i.e.\ $a\%n=r$, where $a=mn+r$.})
\EndFor\EndWhile
\State $X_i=K^{-1}X_i$
\end{algorithmic}\end{algorithm}
